\appendix

\chapter{Appendix}

\section{Full example from Sec. \ref{sec:thread_process}}\label{app:final_example}

\begin{Verbatim}[label=\makebox{\url{https://github.com/lucabaldini/cmepda/tree/master/slides/latex/snippets/benchmark_comparison.py}},commandchars=\\\{\}]
\PY{k}{import} \PY{n+nn}{math}
\PY{k}{import} \PY{n+nn}{multiprocessing}
\PY{k}{import} \PY{n+nn}{random}
\PY{k}{import} \PY{n+nn}{threading}
\PY{k}{import} \PY{n+nn}{time}
\PY{k}{import} \PY{n+nn}{matplotlib}\PY{n+nn}{.}\PY{n+nn}{pyplot} \PY{k}{as} \PY{n+nn}{plt}
\PY{k}{import} \PY{n+nn}{numpy}


\PY{k}{class} \PY{n+nc}{Timer}\PY{p}{(}\PY{n+nb}{object}\PY{p}{)}\PY{p}{:}
    \PY{k}{def} \PY{n+nf+fm}{\_\_init\_\_}\PY{p}{(}\PY{n+nb+bp}{self}\PY{p}{,} \PY{n}{name}\PY{o}{=}\PY{k+kc}{None}\PY{p}{)}\PY{p}{:}
        \PY{n+nb+bp}{self}\PY{o}{.}\PY{n}{name} \PY{o}{=} \PY{n}{name}
        \PY{n+nb+bp}{self}\PY{o}{.}\PY{n}{timee}\PY{o}{=}\PY{l+m+mi}{0}

    \PY{k}{def} \PY{n+nf+fm}{\_\_enter\_\_}\PY{p}{(}\PY{n+nb+bp}{self}\PY{p}{)}\PY{p}{:}
        \PY{n+nb+bp}{self}\PY{o}{.}\PY{n}{tstart} \PY{o}{=} \PY{n}{time}\PY{o}{.}\PY{n}{time}\PY{p}{(}\PY{p}{)}

    \PY{k}{def} \PY{n+nf+fm}{\_\_exit\_\_}\PY{p}{(}\PY{n+nb+bp}{self}\PY{p}{,} \PY{n+nb}{type}\PY{p}{,} \PY{n}{value}\PY{p}{,} \PY{n}{traceback}\PY{p}{)}\PY{p}{:}
        \PY{k}{if} \PY{n+nb+bp}{self}\PY{o}{.}\PY{n}{name}\PY{p}{:}
            \PY{n+nb}{print}\PY{p}{(}\PY{l+s+s1}{\PYZsq{}}\PY{l+s+s1}{[}\PY{l+s+si}{\PYZpc{}s}\PY{l+s+s1}{]}\PY{l+s+s1}{\PYZsq{}} \PY{o}{\PYZpc{}} \PY{n+nb+bp}{self}\PY{o}{.}\PY{n}{name}\PY{p}{,} \PY{n}{end}\PY{o}{=}\PY{l+s+s1}{\PYZsq{}}\PY{l+s+s1}{ }\PY{l+s+s1}{\PYZsq{}}\PY{p}{)}
        \PY{n+nb+bp}{self}\PY{o}{.}\PY{n}{timee}\PY{o}{=}\PY{p}{(}\PY{n}{time}\PY{o}{.}\PY{n}{time}\PY{p}{(}\PY{p}{)} \PY{o}{\PYZhy{}} \PY{n+nb+bp}{self}\PY{o}{.}\PY{n}{tstart}\PY{p}{)}
        \PY{n+nb}{print}\PY{p}{(}\PY{l+s+s1}{\PYZsq{}}\PY{l+s+s1}{Elapsed: }\PY{l+s+si}{\PYZpc{}s}\PY{l+s+s1}{\PYZsq{}} \PY{o}{\PYZpc{}} \PY{p}{(}\PY{n}{time}\PY{o}{.}\PY{n}{time}\PY{p}{(}\PY{p}{)} \PY{o}{\PYZhy{}} \PY{n+nb+bp}{self}\PY{o}{.}\PY{n}{tstart}\PY{p}{)}\PY{p}{)}
        \PY{n+nb+bp}{self}\PY{o}{.}\PY{n}{output}\PY{p}{(}\PY{p}{)}

    \PY{k}{def} \PY{n+nf}{output}\PY{p}{(}\PY{n+nb+bp}{self}\PY{p}{)}\PY{p}{:}
        \PY{k}{return} \PY{n+nb+bp}{self}\PY{o}{.}\PY{n}{timee}


\PY{k}{def} \PY{n+nf}{factorize\_naive}\PY{p}{(}\PY{n}{n}\PY{p}{)}\PY{p}{:}
    \PY{l+s+sd}{\PYZdq{}\PYZdq{}\PYZdq{} A naive factorization method. Take integer \PYZsq{}n\PYZsq{}, return list of}
\PY{l+s+sd}{        factors.}
\PY{l+s+sd}{    \PYZdq{}\PYZdq{}\PYZdq{}}
    \PY{k}{if} \PY{n}{n} \PY{o}{\PYZlt{}} \PY{l+m+mi}{2}\PY{p}{:}
        \PY{k}{return} \PY{p}{[}\PY{p}{]}
    \PY{n}{factors} \PY{o}{=} \PY{p}{[}\PY{p}{]}
    \PY{n}{p} \PY{o}{=} \PY{l+m+mi}{2}

    \PY{k}{while} \PY{k+kc}{True}\PY{p}{:}
        \PY{k}{if} \PY{n}{n} \PY{o}{==} \PY{l+m+mi}{1}\PY{p}{:}
            \PY{k}{return} \PY{n}{factors}
        \PY{n}{r} \PY{o}{=} \PY{n}{n} \PY{o}{\PYZpc{}} \PY{n}{p}
        \PY{k}{if} \PY{n}{r} \PY{o}{==} \PY{l+m+mi}{0}\PY{p}{:}
            \PY{n}{factors}\PY{o}{.}\PY{n}{append}\PY{p}{(}\PY{n}{p}\PY{p}{)}
            \PY{n}{n} \PY{o}{=} \PY{n}{n} \PY{o}{/}\PY{o}{/} \PY{n}{p}
        \PY{k}{elif} \PY{n}{p} \PY{o}{*} \PY{n}{p} \PY{o}{\PYZgt{}}\PY{o}{=} \PY{n}{n}\PY{p}{:}
            \PY{n}{factors}\PY{o}{.}\PY{n}{append}\PY{p}{(}\PY{n}{n}\PY{p}{)}
            \PY{k}{return} \PY{n}{factors}
        \PY{k}{elif} \PY{n}{p} \PY{o}{\PYZgt{}} \PY{l+m+mi}{2}\PY{p}{:}
            \PY{c+c1}{\# Advance in steps of 2 over odd numbers}
            \PY{n}{p} \PY{o}{+}\PY{o}{=} \PY{l+m+mi}{2}
        \PY{k}{else}\PY{p}{:}
            \PY{c+c1}{\# If p == 2, get to 3}
            \PY{n}{p} \PY{o}{+}\PY{o}{=} \PY{l+m+mi}{1}
    \PY{k}{assert} \PY{k+kc}{False}\PY{p}{,} \PY{l+s+s2}{\PYZdq{}}\PY{l+s+s2}{unreachable}\PY{l+s+s2}{\PYZdq{}}


\PY{c+c1}{\# Each \PYZdq{}factorizer\PYZdq{} function returns a dict mapping num \PYZhy{}\PYZgt{} factors}
\PY{k}{def} \PY{n+nf}{serial\_factorizer}\PY{p}{(}\PY{n}{nums}\PY{p}{)}\PY{p}{:}
    \PY{k}{return} \PY{p}{\PYZob{}}\PY{n}{n}\PY{p}{:} \PY{n}{factorize\_naive}\PY{p}{(}\PY{n}{n}\PY{p}{)} \PY{k}{for} \PY{n}{n} \PY{o+ow}{in} \PY{n}{nums}\PY{p}{\PYZcb{}}

\PY{k}{def} \PY{n+nf}{threaded\_factorizer}\PY{p}{(}\PY{n}{nums}\PY{p}{,} \PY{n}{nthreads}\PY{p}{)}\PY{p}{:}
    \PY{k}{def} \PY{n+nf}{worker}\PY{p}{(}\PY{n}{nums}\PY{p}{,} \PY{n}{outdict}\PY{p}{)}\PY{p}{:}
        \PY{l+s+sd}{\PYZdq{}\PYZdq{}\PYZdq{} The worker function, invoked in a thread. \PYZsq{}nums\PYZsq{} is a}
\PY{l+s+sd}{            list of numbers to factor. The results are placed in}
\PY{l+s+sd}{            outdict.}
\PY{l+s+sd}{        \PYZdq{}\PYZdq{}\PYZdq{}}
        \PY{k}{for} \PY{n}{n} \PY{o+ow}{in} \PY{n}{nums}\PY{p}{:}
            \PY{n}{outdict}\PY{p}{[}\PY{n}{n}\PY{p}{]} \PY{o}{=} \PY{n}{factorize\_naive}\PY{p}{(}\PY{n}{n}\PY{p}{)}

    \PY{c+c1}{\# Each thread will get \PYZsq{}chunksize\PYZsq{} nums and its own output dict}
    \PY{n}{chunksize} \PY{o}{=} \PY{n+nb}{int}\PY{p}{(}\PY{n}{math}\PY{o}{.}\PY{n}{ceil}\PY{p}{(}\PY{n+nb}{len}\PY{p}{(}\PY{n}{nums}\PY{p}{)} \PY{o}{/} \PY{n+nb}{float}\PY{p}{(}\PY{n}{nthreads}\PY{p}{)}\PY{p}{)}\PY{p}{)}
    \PY{n}{threads} \PY{o}{=} \PY{p}{[}\PY{p}{]}
    \PY{n}{outs} \PY{o}{=} \PY{p}{[}\PY{p}{\PYZob{}}\PY{p}{\PYZcb{}} \PY{k}{for} \PY{n}{i} \PY{o+ow}{in} \PY{n+nb}{range}\PY{p}{(}\PY{n}{nthreads}\PY{p}{)}\PY{p}{]}

    \PY{k}{for} \PY{n}{i} \PY{o+ow}{in} \PY{n+nb}{range}\PY{p}{(}\PY{n}{nthreads}\PY{p}{)}\PY{p}{:}
        \PY{c+c1}{\# Create each thread, passing it its chunk of numbers to factor}
        \PY{c+c1}{\# and output dict.}
        \PY{n}{t} \PY{o}{=} \PY{n}{threading}\PY{o}{.}\PY{n}{Thread}\PY{p}{(}
                \PY{n}{target}\PY{o}{=}\PY{n}{worker}\PY{p}{,}
                \PY{n}{args}\PY{o}{=}\PY{p}{(}\PY{n}{nums}\PY{p}{[}\PY{n}{chunksize} \PY{o}{*} \PY{n}{i}\PY{p}{:}\PY{n}{chunksize} \PY{o}{*} \PY{p}{(}\PY{n}{i} \PY{o}{+} \PY{l+m+mi}{1}\PY{p}{)}\PY{p}{]}\PY{p}{,}
                      \PY{n}{outs}\PY{p}{[}\PY{n}{i}\PY{p}{]}\PY{p}{)}\PY{p}{)}
        \PY{n}{threads}\PY{o}{.}\PY{n}{append}\PY{p}{(}\PY{n}{t}\PY{p}{)}
        \PY{n}{t}\PY{o}{.}\PY{n}{start}\PY{p}{(}\PY{p}{)}

    \PY{c+c1}{\# Wait for all threads to finish}
    \PY{k}{for} \PY{n}{t} \PY{o+ow}{in} \PY{n}{threads}\PY{p}{:}
        \PY{n}{t}\PY{o}{.}\PY{n}{join}\PY{p}{(}\PY{p}{)}

    \PY{c+c1}{\# Merge all partial output dicts into a single dict and return it}
    \PY{k}{return} \PY{p}{\PYZob{}}\PY{n}{k}\PY{p}{:} \PY{n}{v} \PY{k}{for} \PY{n}{out\_d} \PY{o+ow}{in} \PY{n}{outs} \PY{k}{for} \PY{n}{k}\PY{p}{,} \PY{n}{v} \PY{o+ow}{in} \PY{n}{out\_d}\PY{o}{.}\PY{n}{items}\PY{p}{(}\PY{p}{)}\PY{p}{\PYZcb{}}

\PY{k}{def} \PY{n+nf}{mp\_worker}\PY{p}{(}\PY{n}{nums}\PY{p}{,} \PY{n}{out\_q}\PY{p}{)}\PY{p}{:}
    \PY{l+s+sd}{\PYZdq{}\PYZdq{}\PYZdq{} The worker function, invoked in a process. \PYZsq{}nums\PYZsq{} is a}
\PY{l+s+sd}{        list of numbers to factor. The results are placed in}
\PY{l+s+sd}{        a dictionary that\PYZsq{}s pushed to a queue.}
\PY{l+s+sd}{    \PYZdq{}\PYZdq{}\PYZdq{}}
    \PY{n}{outdict} \PY{o}{=} \PY{p}{\PYZob{}}\PY{p}{\PYZcb{}}
    \PY{k}{for} \PY{n}{n} \PY{o+ow}{in} \PY{n}{nums}\PY{p}{:}
        \PY{n}{outdict}\PY{p}{[}\PY{n}{n}\PY{p}{]} \PY{o}{=} \PY{n}{factorize\_naive}\PY{p}{(}\PY{n}{n}\PY{p}{)}
    \PY{n}{out\_q}\PY{o}{.}\PY{n}{put}\PY{p}{(}\PY{n}{outdict}\PY{p}{)}


\PY{k}{def} \PY{n+nf}{mp\_factorizer}\PY{p}{(}\PY{n}{nums}\PY{p}{,} \PY{n}{nprocs}\PY{p}{)}\PY{p}{:}
    \PY{c+c1}{\# Each process will get \PYZsq{}chunksize\PYZsq{} nums and a queue to put his out}
    \PY{c+c1}{\# dict into}
    \PY{n}{out\_q} \PY{o}{=} \PY{n}{multiprocessing}\PY{o}{.}\PY{n}{Queue}\PY{p}{(}\PY{p}{)}
    \PY{n}{chunksize} \PY{o}{=} \PY{n+nb}{int}\PY{p}{(}\PY{n}{math}\PY{o}{.}\PY{n}{ceil}\PY{p}{(}\PY{n+nb}{len}\PY{p}{(}\PY{n}{nums}\PY{p}{)} \PY{o}{/} \PY{n+nb}{float}\PY{p}{(}\PY{n}{nprocs}\PY{p}{)}\PY{p}{)}\PY{p}{)}
    \PY{n}{procs} \PY{o}{=} \PY{p}{[}\PY{p}{]}

    \PY{k}{for} \PY{n}{i} \PY{o+ow}{in} \PY{n+nb}{range}\PY{p}{(}\PY{n}{nprocs}\PY{p}{)}\PY{p}{:}
        \PY{n}{p} \PY{o}{=} \PY{n}{multiprocessing}\PY{o}{.}\PY{n}{Process}\PY{p}{(}
                \PY{n}{target}\PY{o}{=}\PY{n}{mp\_worker}\PY{p}{,}
                \PY{n}{args}\PY{o}{=}\PY{p}{(}\PY{n}{nums}\PY{p}{[}\PY{n}{chunksize} \PY{o}{*} \PY{n}{i}\PY{p}{:}\PY{n}{chunksize} \PY{o}{*} \PY{p}{(}\PY{n}{i} \PY{o}{+} \PY{l+m+mi}{1}\PY{p}{)}\PY{p}{]}\PY{p}{,}
                      \PY{n}{out\_q}\PY{p}{)}\PY{p}{)}
        \PY{n}{procs}\PY{o}{.}\PY{n}{append}\PY{p}{(}\PY{n}{p}\PY{p}{)}
        \PY{n}{p}\PY{o}{.}\PY{n}{start}\PY{p}{(}\PY{p}{)}

    \PY{c+c1}{\# Collect all results into a single result dict. We know how many dicts}
    \PY{c+c1}{\# with results to expect.}
    \PY{n}{resultdict} \PY{o}{=} \PY{p}{\PYZob{}}\PY{p}{\PYZcb{}}
    \PY{k}{for} \PY{n}{i} \PY{o+ow}{in} \PY{n+nb}{range}\PY{p}{(}\PY{n}{nprocs}\PY{p}{)}\PY{p}{:}
        \PY{n}{resultdict}\PY{o}{.}\PY{n}{update}\PY{p}{(}\PY{n}{out\_q}\PY{o}{.}\PY{n}{get}\PY{p}{(}\PY{p}{)}\PY{p}{)}

    \PY{c+c1}{\# Wait for all worker processes to finish}
    \PY{k}{for} \PY{n}{p} \PY{o+ow}{in} \PY{n}{procs}\PY{p}{:}
        \PY{n}{p}\PY{o}{.}\PY{n}{join}\PY{p}{(}\PY{p}{)}

    \PY{k}{return} \PY{n}{resultdict}

\PY{k}{def} \PY{n+nf}{plot\_results}\PY{p}{(}\PY{n}{elapsed}\PY{p}{)}\PY{p}{:}
    \PY{n}{plt}\PY{o}{.}\PY{n}{rcdefaults}\PY{p}{(}\PY{p}{)}
    \PY{n}{fig}\PY{p}{,} \PY{n}{ax} \PY{o}{=} \PY{n}{plt}\PY{o}{.}\PY{n}{subplots}\PY{p}{(}\PY{p}{)}
    \PY{n}{laby} \PY{o}{=} \PY{p}{(}\PY{l+s+s1}{\PYZsq{}}\PY{l+s+s1}{Serial}\PY{l+s+s1}{\PYZsq{}}\PY{p}{,}\PY{l+s+s1}{\PYZsq{}}\PY{l+s+s1}{Thread 2}\PY{l+s+s1}{\PYZsq{}}\PY{p}{,}\PY{l+s+s1}{\PYZsq{}}\PY{l+s+s1}{Process 2}\PY{l+s+s1}{\PYZsq{}}\PY{p}{,}\PY{l+s+s1}{\PYZsq{}}\PY{l+s+s1}{Thread 4}\PY{l+s+s1}{\PYZsq{}}\PY{p}{,}\PY{l+s+s1}{\PYZsq{}}\PY{l+s+s1}{Process 4}\PY{l+s+s1}{\PYZsq{}}\PY{p}{,}\PY{l+s+s1}{\PYZsq{}}\PY{l+s+s1}{Thread 8}\PY{l+s+s1}{\PYZsq{}}\PY{p}{,}\PY{l+s+s1}{\PYZsq{}}\PY{l+s+s1}{Process 8}\PY{l+s+s1}{\PYZsq{}}\PY{p}{)}
    \PY{n}{y\_pos} \PY{o}{=} \PY{n}{numpy}\PY{o}{.}\PY{n}{arange}\PY{p}{(}\PY{n+nb}{len}\PY{p}{(}\PY{n}{laby}\PY{p}{)}\PY{p}{)}
    \PY{n}{ax}\PY{o}{.}\PY{n}{barh}\PY{p}{(}\PY{n}{y\_pos}\PY{p}{,} \PY{n}{elapsed}\PY{p}{,} \PY{n}{align}\PY{o}{=}\PY{l+s+s1}{\PYZsq{}}\PY{l+s+s1}{center}\PY{l+s+s1}{\PYZsq{}}\PY{p}{)}
    \PY{n}{ax}\PY{o}{.}\PY{n}{set\_yticks}\PY{p}{(}\PY{n}{y\_pos}\PY{p}{)}
    \PY{n}{ax}\PY{o}{.}\PY{n}{set\_yticklabels}\PY{p}{(}\PY{n}{laby}\PY{p}{)}
    \PY{n}{ax}\PY{o}{.}\PY{n}{invert\_yaxis}\PY{p}{(}\PY{p}{)}  \PY{c+c1}{\# labels read top\PYZhy{}to\PYZhy{}bottom}
    \PY{n}{ax}\PY{o}{.}\PY{n}{set\_xlabel}\PY{p}{(}\PY{l+s+s1}{\PYZsq{}}\PY{l+s+s1}{Elapsed time [s]}\PY{l+s+s1}{\PYZsq{}}\PY{p}{)}
    \PY{n}{ax}\PY{o}{.}\PY{n}{set\_title}\PY{p}{(}\PY{l+s+s1}{\PYZsq{}}\PY{l+s+s1}{Serial, threads, processes comparison}\PY{l+s+s1}{\PYZsq{}}\PY{p}{)}
    \PY{n}{plt}\PY{o}{.}\PY{n}{show}\PY{p}{(}\PY{p}{)}

\PY{k}{def} \PY{n+nf}{benchmark}\PY{p}{(}\PY{n}{nums}\PY{p}{)}\PY{p}{:}
    \PY{n+nb}{print}\PY{p}{(}\PY{l+s+s1}{\PYZsq{}}\PY{l+s+s1}{Running benchmark...}\PY{l+s+s1}{\PYZsq{}}\PY{p}{)}
    \PY{n}{elapsed\_times}\PY{o}{=}\PY{p}{[}\PY{p}{]}

    \PY{n}{tserial}\PY{o}{=}\PY{n}{Timer}\PY{p}{(}\PY{l+s+s1}{\PYZsq{}}\PY{l+s+s1}{serial}\PY{l+s+s1}{\PYZsq{}}\PY{p}{)}
    \PY{k}{with} \PY{n}{tserial} \PY{k}{as} \PY{n}{qq}\PY{p}{:}
        \PY{n}{s\_d} \PY{o}{=} \PY{n}{serial\_factorizer}\PY{p}{(}\PY{n}{nums}\PY{p}{)}
    \PY{n}{elapsed\_times}\PY{o}{.}\PY{n}{append}\PY{p}{(}\PY{n}{tserial}\PY{o}{.}\PY{n}{output}\PY{p}{(}\PY{p}{)}\PY{p}{)}

    \PY{k}{for} \PY{n}{numparallel} \PY{o+ow}{in} \PY{p}{[}\PY{l+m+mi}{2}\PY{p}{,} \PY{l+m+mi}{4}\PY{p}{,} \PY{l+m+mi}{8}\PY{p}{]}\PY{p}{:}
        \PY{n}{tthread}\PY{o}{=}\PY{n}{Timer}\PY{p}{(}\PY{l+s+s1}{\PYZsq{}}\PY{l+s+s1}{threaded }\PY{l+s+si}{\PYZpc{}s}\PY{l+s+s1}{\PYZsq{}} \PY{o}{\PYZpc{}} \PY{n}{numparallel}\PY{p}{)}
        \PY{k}{with} \PY{n}{tthread} \PY{k}{as} \PY{n}{qq}\PY{p}{:}
            \PY{n}{t\_d} \PY{o}{=} \PY{n}{threaded\_factorizer}\PY{p}{(}\PY{n}{nums}\PY{p}{,} \PY{n}{numparallel}\PY{p}{)}
        \PY{n}{elapsed\_times}\PY{o}{.}\PY{n}{append}\PY{p}{(}\PY{n}{tthread}\PY{o}{.}\PY{n}{output}\PY{p}{(}\PY{p}{)}\PY{p}{)}
        \PY{n}{tmpar}\PY{o}{=}\PY{n}{Timer}\PY{p}{(}\PY{l+s+s1}{\PYZsq{}}\PY{l+s+s1}{mp }\PY{l+s+si}{\PYZpc{}s}\PY{l+s+s1}{\PYZsq{}} \PY{o}{\PYZpc{}} \PY{n}{numparallel}\PY{p}{)}
        \PY{k}{with} \PY{n}{tmpar} \PY{k}{as} \PY{n}{qq}\PY{p}{:}
            \PY{n}{m\_d} \PY{o}{=} \PY{n}{mp\_factorizer}\PY{p}{(}\PY{n}{nums}\PY{p}{,} \PY{n}{numparallel}\PY{p}{)}
        \PY{n}{elapsed\_times}\PY{o}{.}\PY{n}{append}\PY{p}{(}\PY{n}{tmpar}\PY{o}{.}\PY{n}{output}\PY{p}{(}\PY{p}{)}\PY{p}{)}

    \PY{n+nb}{print} \PY{p}{(}\PY{n}{elapsed\_times}\PY{p}{)}
    \PY{n}{plot\_results}\PY{p}{(}\PY{n}{elapsed\_times}\PY{p}{)}


\PY{c+c1}{\#MAIN}
\PY{k}{if} \PY{n}{\_\_name\_\_} \PY{o}{==} \PY{l+s+s1}{\PYZsq{}}\PY{l+s+s1}{\_\_main\_\_}\PY{l+s+s1}{\PYZsq{}}\PY{p}{:}
    \PY{n}{N} \PY{o}{=} \PY{l+m+mi}{299}

    \PY{n}{nums} \PY{o}{=} \PY{p}{[}\PY{l+m+mi}{999999999999}\PY{p}{]}
    \PY{k}{for} \PY{n}{i} \PY{o+ow}{in} \PY{n+nb}{range}\PY{p}{(}\PY{n}{N}\PY{p}{)}\PY{p}{:}
        \PY{n}{nums}\PY{o}{.}\PY{n}{append}\PY{p}{(}\PY{n}{nums}\PY{p}{[}\PY{o}{\PYZhy{}}\PY{l+m+mi}{1}\PY{p}{]} \PY{o}{+} \PY{l+m+mi}{2}\PY{p}{)}

    \PY{n+nb}{print}\PY{p}{(}\PY{n}{nums}\PY{p}{)}
    \PY{n}{benchmark}\PY{p}{(}\PY{n}{nums}\PY{p}{)}
\end{Verbatim}

\section{C problems and exercises}\label{app:C_exercises}

\section{Full example from Sec. \ref{sec:hands_on_cuda_1}}\label{app:matrix_example_cuda}

File: \underline{MatrixMultiplication\_global.h}

\begin{Verbatim}[commandchars=\\\{\}]
\PY{c+cp}{\#}\PY{c+cp}{include}\PY{+w}{ }\PY{c+cpf}{\PYZlt{}stdio.h\PYZgt{}}

\PY{c+c1}{// Matrices are stored in row\PYZhy{}major order:}
\PY{c+c1}{// M(row, col) = *(M.elements + row * M.width + col)}
\PY{k}{typedef}\PY{+w}{ }\PY{k}{struct}\PY{+w}{ }\PY{p}{\PYZob{}}
\PY{+w}{    }\PY{k+kt}{int}\PY{+w}{ }\PY{n}{width}\PY{p}{;}
\PY{+w}{    }\PY{k+kt}{int}\PY{+w}{ }\PY{n}{height}\PY{p}{;}
\PY{+w}{    }\PY{k+kt}{float}\PY{o}{*}\PY{+w}{ }\PY{n}{elements}\PY{p}{;}
\PY{p}{\PYZcb{}}\PY{+w}{ }\PY{n}{Matrix}\PY{p}{;}

\PY{c+c1}{// Thread block size}
\PY{c+cp}{\#}\PY{c+cp}{define BLOCK\_SIZE 16}

\PY{k+kr}{\_\_global\_\_}\PY{+w}{ }\PY{k+kt}{void}\PY{+w}{ }\PY{n}{MatMulKernel}\PY{p}{(}\PY{k}{const}\PY{+w}{ }\PY{n}{Matrix}\PY{p}{,}\PY{+w}{ }\PY{k}{const}\PY{+w}{ }\PY{n}{Matrix}\PY{p}{,}\PY{+w}{ }\PY{n}{Matrix}\PY{p}{)}\PY{p}{;}
\end{Verbatim}

\hrulefill

File: \underline{MatrixMultiplication\_global.cu}. \\ Note that when you run the corresponding executable you should specify
\begin{itemize}
    \item[-] the height of the first matrix;
    \item[-] the length of the first matrix (which is set by default equal to the height of the second matrix);
    \item[-] the length of the second matrix.
\end{itemize}
In our case we run \texttt{MatrixMultiplication\_global 1024 1024 1024}.

\begin{Verbatim}[label=\makebox{\url{https://github.com/lucabaldini/cmepda/tree/master/slides/latex/snippets/MatrixMultiplication_global.cu}},commandchars=\\\{\}]
\PY{c+cp}{\#}\PY{c+cp}{include}\PY{+w}{ }\PY{c+cpf}{\PYZlt{}stdlib.h\PYZgt{}}
\PY{c+cp}{\#}\PY{c+cp}{include}\PY{+w}{ }\PY{c+cpf}{\PYZdq{}MatrixMultiplication\_global.h\PYZdq{}}

\PY{c+c1}{// Matrix multiplication \PYZhy{} Host code}
\PY{c+c1}{// Matrix dimensions are assumed to be multiples of BLOCK\_SIZE}
\PY{k+kt}{void}\PY{+w}{ }\PY{n+nf}{MatMul}\PY{p}{(}\PY{k}{const}\PY{+w}{ }\PY{n}{Matrix}\PY{+w}{ }\PY{n}{A}\PY{p}{,}\PY{+w}{ }\PY{k}{const}\PY{+w}{ }\PY{n}{Matrix}\PY{+w}{ }\PY{n}{B}\PY{p}{,}\PY{+w}{ }\PY{n}{Matrix}\PY{+w}{ }\PY{n}{C}\PY{p}{)}\PY{+w}{ }\PY{p}{\PYZob{}}

\PY{+w}{  }\PY{c+c1}{// \PYZhy{}\PYZhy{}\PYZhy{} 1. ALLOCATE AND COPY MATRIX A TO DEVICE \PYZhy{}\PYZhy{}\PYZhy{}}
\PY{+w}{  }\PY{n}{Matrix}\PY{+w}{ }\PY{n}{d\_A}\PY{p}{;}
\PY{+w}{  }\PY{n}{d\_A}\PY{p}{.}\PY{n}{width}\PY{+w}{ }\PY{o}{=}\PY{+w}{ }\PY{n}{A}\PY{p}{.}\PY{n}{width}\PY{p}{;}
\PY{+w}{  }\PY{n}{d\_A}\PY{p}{.}\PY{n}{height}\PY{+w}{ }\PY{o}{=}\PY{+w}{ }\PY{n}{A}\PY{p}{.}\PY{n}{height}\PY{p}{;}

\PY{+w}{  }\PY{c+c1}{// Calculate the total size in bytes for the 1D array representing Matrix A}
\PY{+w}{  }\PY{k+kt}{size\_t}\PY{+w}{ }\PY{n}{size}\PY{+w}{ }\PY{o}{=}\PY{+w}{ }\PY{n}{A}\PY{p}{.}\PY{n}{width}\PY{+w}{ }\PY{o}{*}\PY{+w}{ }\PY{n}{A}\PY{p}{.}\PY{n}{height}\PY{+w}{ }\PY{o}{*}\PY{+w}{ }\PY{k}{sizeof}\PY{p}{(}\PY{k+kt}{float}\PY{p}{)}\PY{p}{;}

\PY{+w}{  }\PY{c+c1}{// Allocate memory on the GPU (Device)}
\PY{+w}{  }\PY{n}{cudaError\_t}\PY{+w}{ }\PY{n}{err}\PY{+w}{ }\PY{o}{=}\PY{+w}{ }\PY{n}{cudaMalloc}\PY{p}{(}\PY{o}{\PYZam{}}\PY{n}{d\_A}\PY{p}{.}\PY{n}{elements}\PY{p}{,}\PY{+w}{ }\PY{n}{size}\PY{p}{)}\PY{p}{;}
\PY{+w}{  }\PY{n}{printf}\PY{p}{(}\PY{l+s}{\PYZdq{}}\PY{l+s}{CUDA malloc A: \PYZpc{}s}\PY{l+s+se}{\PYZbs{}n}\PY{l+s}{\PYZdq{}}\PY{p}{,}\PY{n}{cudaGetErrorString}\PY{p}{(}\PY{n}{err}\PY{p}{)}\PY{p}{)}\PY{p}{;}

\PY{+w}{  }\PY{c+c1}{// Copy data from Host (CPU) to Device (GPU)}
\PY{+w}{  }\PY{n}{err}\PY{+w}{ }\PY{o}{=}\PY{+w}{ }\PY{n}{cudaMemcpy}\PY{p}{(}\PY{n}{d\_A}\PY{p}{.}\PY{n}{elements}\PY{p}{,}\PY{+w}{ }\PY{n}{A}\PY{p}{.}\PY{n}{elements}\PY{p}{,}\PY{+w}{ }\PY{n}{size}\PY{p}{,}\PY{+w}{ }\PY{n}{cudaMemcpyHostToDevice}\PY{p}{)}\PY{p}{;}
\PY{+w}{  }\PY{n}{printf}\PY{p}{(}\PY{l+s}{\PYZdq{}}\PY{l+s}{Copy A to device: \PYZpc{}s}\PY{l+s+se}{\PYZbs{}n}\PY{l+s}{\PYZdq{}}\PY{p}{,}\PY{n}{cudaGetErrorString}\PY{p}{(}\PY{n}{err}\PY{p}{)}\PY{p}{)}\PY{p}{;}

\PY{+w}{  }\PY{c+c1}{// \PYZhy{}\PYZhy{}\PYZhy{} 2. ALLOCATE AND COPY MATRIX B TO DEVICE \PYZhy{}\PYZhy{}\PYZhy{}}
\PY{+w}{  }\PY{n}{Matrix}\PY{+w}{ }\PY{n}{d\_B}\PY{p}{;}
\PY{+w}{  }\PY{n}{d\_B}\PY{p}{.}\PY{n}{width}\PY{+w}{ }\PY{o}{=}\PY{+w}{ }\PY{n}{B}\PY{p}{.}\PY{n}{width}\PY{p}{;}
\PY{+w}{  }\PY{n}{d\_B}\PY{p}{.}\PY{n}{height}\PY{+w}{ }\PY{o}{=}\PY{+w}{ }\PY{n}{B}\PY{p}{.}\PY{n}{height}\PY{p}{;}
\PY{+w}{  }\PY{n}{size}\PY{+w}{ }\PY{o}{=}\PY{+w}{ }\PY{n}{B}\PY{p}{.}\PY{n}{width}\PY{+w}{ }\PY{o}{*}\PY{+w}{ }\PY{n}{B}\PY{p}{.}\PY{n}{height}\PY{+w}{ }\PY{o}{*}\PY{+w}{ }\PY{k}{sizeof}\PY{p}{(}\PY{k+kt}{float}\PY{p}{)}\PY{p}{;}

\PY{+w}{  }\PY{n}{err}\PY{+w}{ }\PY{o}{=}\PY{+w}{ }\PY{n}{cudaMalloc}\PY{p}{(}\PY{o}{\PYZam{}}\PY{n}{d\_B}\PY{p}{.}\PY{n}{elements}\PY{p}{,}\PY{+w}{ }\PY{n}{size}\PY{p}{)}\PY{p}{;}
\PY{+w}{  }\PY{n}{printf}\PY{p}{(}\PY{l+s}{\PYZdq{}}\PY{l+s}{CUDA malloc B: \PYZpc{}s}\PY{l+s+se}{\PYZbs{}n}\PY{l+s}{\PYZdq{}}\PY{p}{,}\PY{n}{cudaGetErrorString}\PY{p}{(}\PY{n}{err}\PY{p}{)}\PY{p}{)}\PY{p}{;}
\PY{+w}{  }\PY{n}{err}\PY{+w}{ }\PY{o}{=}\PY{+w}{ }\PY{n}{cudaMemcpy}\PY{p}{(}\PY{n}{d\_B}\PY{p}{.}\PY{n}{elements}\PY{p}{,}\PY{+w}{ }\PY{n}{B}\PY{p}{.}\PY{n}{elements}\PY{p}{,}\PY{+w}{ }\PY{n}{size}\PY{p}{,}\PY{+w}{ }\PY{n}{cudaMemcpyHostToDevice}\PY{p}{)}\PY{p}{;}
\PY{+w}{  }\PY{n}{printf}\PY{p}{(}\PY{l+s}{\PYZdq{}}\PY{l+s}{Copy B to device: \PYZpc{}s}\PY{l+s+se}{\PYZbs{}n}\PY{l+s}{\PYZdq{}}\PY{p}{,}\PY{n}{cudaGetErrorString}\PY{p}{(}\PY{n}{err}\PY{p}{)}\PY{p}{)}\PY{p}{;}

\PY{+w}{  }\PY{c+c1}{// \PYZhy{}\PYZhy{}\PYZhy{} 3. ALLOCATE MATRIX C ON DEVICE \PYZhy{}\PYZhy{}\PYZhy{}}
\PY{+w}{  }\PY{n}{Matrix}\PY{+w}{ }\PY{n}{d\_C}\PY{p}{;}
\PY{+w}{  }\PY{n}{d\_C}\PY{p}{.}\PY{n}{width}\PY{+w}{ }\PY{o}{=}\PY{+w}{ }\PY{n}{C}\PY{p}{.}\PY{n}{width}\PY{p}{;}
\PY{+w}{  }\PY{n}{d\_C}\PY{p}{.}\PY{n}{height}\PY{+w}{ }\PY{o}{=}\PY{+w}{ }\PY{n}{C}\PY{p}{.}\PY{n}{height}\PY{p}{;}
\PY{+w}{  }\PY{n}{size}\PY{+w}{ }\PY{o}{=}\PY{+w}{ }\PY{n}{C}\PY{p}{.}\PY{n}{width}\PY{+w}{ }\PY{o}{*}\PY{+w}{ }\PY{n}{C}\PY{p}{.}\PY{n}{height}\PY{+w}{ }\PY{o}{*}\PY{+w}{ }\PY{k}{sizeof}\PY{p}{(}\PY{k+kt}{float}\PY{p}{)}\PY{p}{;}

\PY{+w}{  }\PY{n}{err}\PY{+w}{ }\PY{o}{=}\PY{+w}{ }\PY{n}{cudaMalloc}\PY{p}{(}\PY{o}{\PYZam{}}\PY{n}{d\_C}\PY{p}{.}\PY{n}{elements}\PY{p}{,}\PY{+w}{ }\PY{n}{size}\PY{p}{)}\PY{p}{;}
\PY{+w}{  }\PY{n}{printf}\PY{p}{(}\PY{l+s}{\PYZdq{}}\PY{l+s}{CUDA malloc C: \PYZpc{}s}\PY{l+s+se}{\PYZbs{}n}\PY{l+s}{\PYZdq{}}\PY{p}{,}\PY{n}{cudaGetErrorString}\PY{p}{(}\PY{n}{err}\PY{p}{)}\PY{p}{)}\PY{p}{;}

\PY{+w}{  }\PY{c+c1}{// \PYZhy{}\PYZhy{}\PYZhy{} 4. TIMING SETUP \PYZhy{}\PYZhy{}\PYZhy{}}
\PY{+w}{  }\PY{k+kt}{float}\PY{+w}{ }\PY{n}{time}\PY{p}{;}
\PY{+w}{  }\PY{n}{cudaEvent\_t}\PY{+w}{ }\PY{n}{start}\PY{p}{,}\PY{+w}{ }\PY{n}{stop}\PY{p}{;}
\PY{+w}{  }\PY{n}{cudaEventCreate}\PY{p}{(}\PY{o}{\PYZam{}}\PY{n}{start}\PY{p}{)}\PY{p}{;}
\PY{+w}{  }\PY{n}{cudaEventCreate}\PY{p}{(}\PY{o}{\PYZam{}}\PY{n}{stop}\PY{p}{)}\PY{p}{;}

\PY{+w}{  }\PY{c+c1}{// Start the timer}
\PY{+w}{  }\PY{n}{cudaEventRecord}\PY{p}{(}\PY{n}{start}\PY{p}{)}\PY{p}{;}

\PY{+w}{  }\PY{c+c1}{// \PYZhy{}\PYZhy{}\PYZhy{} 5. GRID AND BLOCK CONFIGURATION \PYZhy{}\PYZhy{}\PYZhy{}}
\PY{+w}{  }\PY{c+c1}{// Define the number of threads per block in 2D (BLOCK\_SIZE x BLOCK\_SIZE)}
\PY{+w}{  }\PY{k+kt}{dim3}\PY{+w}{ }\PY{n}{dimBlock}\PY{p}{(}\PY{n}{BLOCK\_SIZE}\PY{p}{,}\PY{+w}{ }\PY{n}{BLOCK\_SIZE}\PY{p}{)}\PY{p}{;}

\PY{+w}{  }\PY{c+c1}{// Calculate the number of blocks needed in the grid (X and Y dimensions)}
\PY{+w}{  }\PY{c+c1}{// Adding (dimBlock \PYZhy{} 1) before dividing ensures we round up if dimensions}
\PY{+w}{  }\PY{c+c1}{// are not perfect multiples of BLOCK\_SIZE.}
\PY{+w}{  }\PY{k+kt}{dim3}\PY{+w}{ }\PY{n}{dimGrid}\PY{p}{(}\PY{p}{(}\PY{n}{B}\PY{p}{.}\PY{n}{width}\PY{+w}{ }\PY{o}{+}\PY{+w}{ }\PY{n}{dimBlock}\PY{p}{.}\PY{n}{x}\PY{+w}{ }\PY{o}{\PYZhy{}}\PY{+w}{ }\PY{l+m+mi}{1}\PY{p}{)}\PY{+w}{ }\PY{o}{/}\PY{+w}{ }\PY{n}{dimBlock}\PY{p}{.}\PY{n}{x}\PY{p}{,}
\PY{+w}{               }\PY{p}{(}\PY{n}{A}\PY{p}{.}\PY{n}{height}\PY{+w}{ }\PY{o}{+}\PY{+w}{ }\PY{n}{dimBlock}\PY{p}{.}\PY{n}{y}\PY{+w}{ }\PY{o}{\PYZhy{}}\PY{+w}{ }\PY{l+m+mi}{1}\PY{p}{)}\PY{+w}{ }\PY{o}{/}\PY{+w}{ }\PY{n}{dimBlock}\PY{p}{.}\PY{n}{y}\PY{p}{)}\PY{p}{;}

\PY{+w}{  }\PY{c+c1}{// Launch the kernel}
\PY{+w}{  }\PY{n}{MatMulKernel}\PY{o}{\PYZlt{}}\PY{o}{\PYZlt{}}\PY{o}{\PYZlt{}}\PY{n}{dimGrid}\PY{p}{,}\PY{+w}{ }\PY{n}{dimBlock}\PY{o}{\PYZgt{}}\PY{o}{\PYZgt{}}\PY{o}{\PYZgt{}}\PY{p}{(}\PY{n}{d\_A}\PY{p}{,}\PY{+w}{ }\PY{n}{d\_B}\PY{p}{,}\PY{+w}{ }\PY{n}{d\_C}\PY{p}{)}\PY{p}{;}

\PY{+w}{  }\PY{c+c1}{// Wait for all threads to finish}
\PY{+w}{  }\PY{n}{cudaDeviceSynchronize}\PY{p}{(}\PY{p}{)}\PY{p}{;}

\PY{+w}{  }\PY{c+c1}{// \PYZhy{}\PYZhy{}\PYZhy{} 6. STOP TIMER AND PRINT \PYZhy{}\PYZhy{}\PYZhy{}}
\PY{+w}{  }\PY{n}{cudaEventRecord}\PY{p}{(}\PY{n}{stop}\PY{p}{)}\PY{p}{;}
\PY{+w}{  }\PY{n}{cudaEventSynchronize}\PY{p}{(}\PY{n}{stop}\PY{p}{)}\PY{p}{;}
\PY{+w}{  }\PY{n}{cudaEventElapsedTime}\PY{p}{(}\PY{o}{\PYZam{}}\PY{n}{time}\PY{p}{,}\PY{+w}{ }\PY{n}{start}\PY{p}{,}\PY{+w}{ }\PY{n}{stop}\PY{p}{)}\PY{p}{;}

\PY{+w}{  }\PY{n}{printf}\PY{p}{(}\PY{l+s}{\PYZdq{}}\PY{l+s}{Run kernel: \PYZpc{}s}\PY{l+s+se}{\PYZbs{}n}\PY{l+s}{\PYZdq{}}\PY{p}{,}\PY{+w}{ }\PY{n}{cudaGetErrorString}\PY{p}{(}\PY{n}{err}\PY{p}{)}\PY{p}{)}\PY{p}{;}
\PY{+w}{  }\PY{n}{printf}\PY{p}{(}\PY{l+s}{\PYZdq{}}\PY{l+s}{Time: \PYZpc{}3.5f ms}\PY{l+s+se}{\PYZbs{}n}\PY{l+s}{\PYZdq{}}\PY{p}{,}\PY{+w}{ }\PY{n}{time}\PY{p}{)}\PY{p}{;}

\PY{+w}{  }\PY{c+c1}{// \PYZhy{}\PYZhy{}\PYZhy{} 7. RETRIEVE RESULTS AND CLEANUP \PYZhy{}\PYZhy{}\PYZhy{}}
\PY{+w}{  }\PY{c+c1}{// Copy the computed matrix C back to the Host}
\PY{+w}{  }\PY{n}{err}\PY{+w}{ }\PY{o}{=}\PY{+w}{ }\PY{n}{cudaMemcpy}\PY{p}{(}\PY{n}{C}\PY{p}{.}\PY{n}{elements}\PY{p}{,}\PY{+w}{ }\PY{n}{d\_C}\PY{p}{.}\PY{n}{elements}\PY{p}{,}\PY{+w}{ }\PY{n}{size}\PY{p}{,}\PY{+w}{ }\PY{n}{cudaMemcpyDeviceToHost}\PY{p}{)}\PY{p}{;}
\PY{+w}{  }\PY{n}{printf}\PY{p}{(}\PY{l+s}{\PYZdq{}}\PY{l+s}{Copy C off of device: \PYZpc{}s}\PY{l+s+se}{\PYZbs{}n}\PY{l+s}{\PYZdq{}}\PY{p}{,}\PY{n}{cudaGetErrorString}\PY{p}{(}\PY{n}{err}\PY{p}{)}\PY{p}{)}\PY{p}{;}

\PY{+w}{  }\PY{c+c1}{// Free the device memory to prevent memory leaks}
\PY{+w}{  }\PY{n}{cudaFree}\PY{p}{(}\PY{n}{d\_A}\PY{p}{.}\PY{n}{elements}\PY{p}{)}\PY{p}{;}
\PY{+w}{  }\PY{n}{cudaFree}\PY{p}{(}\PY{n}{d\_B}\PY{p}{.}\PY{n}{elements}\PY{p}{)}\PY{p}{;}
\PY{+w}{  }\PY{n}{cudaFree}\PY{p}{(}\PY{n}{d\_C}\PY{p}{.}\PY{n}{elements}\PY{p}{)}\PY{p}{;}
\PY{p}{\PYZcb{}}

\PY{c+c1}{// Matrix multiplication kernel called by MatMul()}
\PY{k+kr}{\_\_global\_\_}\PY{+w}{ }\PY{k+kt}{void}\PY{+w}{ }\PY{n}{MatMulKernel}\PY{p}{(}\PY{n}{Matrix}\PY{+w}{ }\PY{n}{A}\PY{p}{,}\PY{+w}{ }\PY{n}{Matrix}\PY{+w}{ }\PY{n}{B}\PY{p}{,}\PY{+w}{ }\PY{n}{Matrix}\PY{+w}{ }\PY{n}{C}\PY{p}{)}\PY{+w}{ }\PY{p}{\PYZob{}}

\PY{+w}{  }\PY{c+c1}{// Variable to accumulate the dot product sum for this specific thread\PYZsq{}s element}
\PY{+w}{  }\PY{k+kt}{float}\PY{+w}{ }\PY{n}{Cvalue}\PY{+w}{ }\PY{o}{=}\PY{+w}{ }\PY{l+m+mf}{0.0}\PY{p}{;}

\PY{+w}{  }\PY{c+c1}{// Calculate the global 2D row and column index for the current thread}
\PY{+w}{  }\PY{c+c1}{// blockIdx: Which block we are in}
\PY{+w}{  }\PY{c+c1}{// blockDim: Size of the block}
\PY{+w}{  }\PY{c+c1}{// threadIdx: Which thread we are inside the block}
\PY{+w}{  }\PY{k+kt}{int}\PY{+w}{ }\PY{n}{row}\PY{+w}{ }\PY{o}{=}\PY{+w}{ }\PY{n+nb}{blockIdx}\PY{p}{.}\PY{n}{y}\PY{+w}{ }\PY{o}{*}\PY{+w}{ }\PY{n+nb}{blockDim}\PY{p}{.}\PY{n}{y}\PY{+w}{ }\PY{o}{+}\PY{+w}{ }\PY{n+nb}{threadIdx}\PY{p}{.}\PY{n}{y}\PY{p}{;}
\PY{+w}{  }\PY{k+kt}{int}\PY{+w}{ }\PY{n}{col}\PY{+w}{ }\PY{o}{=}\PY{+w}{ }\PY{n+nb}{blockIdx}\PY{p}{.}\PY{n}{x}\PY{+w}{ }\PY{o}{*}\PY{+w}{ }\PY{n+nb}{blockDim}\PY{p}{.}\PY{n}{x}\PY{+w}{ }\PY{o}{+}\PY{+w}{ }\PY{n+nb}{threadIdx}\PY{p}{.}\PY{n}{x}\PY{p}{;}

\PY{+w}{  }\PY{c+c1}{// Boundary check: ensure the thread maps to a valid matrix coordinate}
\PY{+w}{  }\PY{k}{if}\PY{p}{(}\PY{n}{row}\PY{+w}{ }\PY{o}{\PYZgt{}}\PY{+w}{ }\PY{n}{A}\PY{p}{.}\PY{n}{height}\PY{+w}{ }\PY{o}{|}\PY{o}{|}\PY{+w}{ }\PY{n}{col}\PY{+w}{ }\PY{o}{\PYZgt{}}\PY{o}{=}\PY{+w}{ }\PY{n}{B}\PY{p}{.}\PY{n}{width}\PY{p}{)}\PY{+w}{ }\PY{k}{return}\PY{p}{;}

\PY{+w}{  }\PY{c+c1}{// Perform the dot product for the calculated row of A and column of B}
\PY{+w}{  }\PY{k}{for}\PY{+w}{ }\PY{p}{(}\PY{k+kt}{int}\PY{+w}{ }\PY{n}{e}\PY{+w}{ }\PY{o}{=}\PY{+w}{ }\PY{l+m+mi}{0}\PY{p}{;}\PY{+w}{ }\PY{n}{e}\PY{+w}{ }\PY{o}{\PYZlt{}}\PY{+w}{ }\PY{n}{A}\PY{p}{.}\PY{n}{width}\PY{p}{;}\PY{+w}{ }\PY{o}{+}\PY{o}{+}\PY{n}{e}\PY{p}{)}\PY{+w}{ }\PY{p}{\PYZob{}}
\PY{+w}{    }\PY{c+c1}{// 2D arrays are flattened to 1D in C.}
\PY{+w}{    }\PY{c+c1}{// Formula to access (row, col) is: [row * width + col]}
\PY{+w}{    }\PY{n}{Cvalue}\PY{+w}{ }\PY{o}{+}\PY{o}{=}\PY{+w}{ }\PY{p}{(}\PY{n}{A}\PY{p}{.}\PY{n}{elements}\PY{p}{[}\PY{n}{row}\PY{+w}{ }\PY{o}{*}\PY{+w}{ }\PY{n}{A}\PY{p}{.}\PY{n}{width}\PY{+w}{ }\PY{o}{+}\PY{+w}{ }\PY{n}{e}\PY{p}{]}\PY{p}{)}\PY{+w}{ }\PY{o}{*}\PY{+w}{ }\PY{p}{(}\PY{n}{B}\PY{p}{.}\PY{n}{elements}\PY{p}{[}\PY{n}{e}\PY{+w}{ }\PY{o}{*}\PY{+w}{ }\PY{n}{B}\PY{p}{.}\PY{n}{width}\PY{+w}{ }\PY{o}{+}\PY{+w}{ }\PY{n}{col}\PY{p}{]}\PY{p}{)}\PY{p}{;}
\PY{+w}{  }\PY{p}{\PYZcb{}}

\PY{+w}{  }\PY{c+c1}{// Write the final accumulated value to the correct 1D index in output Matrix C}
\PY{+w}{  }\PY{n}{C}\PY{p}{.}\PY{n}{elements}\PY{p}{[}\PY{n}{row}\PY{+w}{ }\PY{o}{*}\PY{+w}{ }\PY{n}{C}\PY{p}{.}\PY{n}{width}\PY{+w}{ }\PY{o}{+}\PY{+w}{ }\PY{n}{col}\PY{p}{]}\PY{+w}{ }\PY{o}{=}\PY{+w}{ }\PY{n}{Cvalue}\PY{p}{;}
\PY{p}{\PYZcb{}}

\PY{c+c1}{// Usage: .\PYZbs{}MatrixMultiplication\_global a1 a2 b2}
\PY{k+kt}{int}\PY{+w}{ }\PY{n}{main}\PY{p}{(}\PY{k+kt}{int}\PY{+w}{ }\PY{n}{argc}\PY{p}{,}\PY{+w}{ }\PY{k+kt}{char}\PY{o}{*}\PY{+w}{ }\PY{n}{argv}\PY{p}{[}\PY{p}{]}\PY{p}{)}\PY{p}{\PYZob{}}
\PY{+w}{  }\PY{n}{Matrix}\PY{+w}{ }\PY{n}{A}\PY{p}{,}\PY{+w}{ }\PY{n}{B}\PY{p}{,}\PY{+w}{ }\PY{n}{C}\PY{p}{;}
\PY{+w}{  }\PY{k+kt}{int}\PY{+w}{ }\PY{n}{a1}\PY{p}{,}\PY{+w}{ }\PY{n}{a2}\PY{p}{,}\PY{+w}{ }\PY{n}{b1}\PY{p}{,}\PY{+w}{ }\PY{n}{b2}\PY{p}{;}

\PY{+w}{  }\PY{c+c1}{// Read matrix dimensions from command line arguments}
\PY{+w}{  }\PY{n}{a1}\PY{+w}{ }\PY{o}{=}\PY{+w}{ }\PY{n}{atoi}\PY{p}{(}\PY{n}{argv}\PY{p}{[}\PY{l+m+mi}{1}\PY{p}{]}\PY{p}{)}\PY{p}{;}\PY{+w}{ }\PY{c+c1}{// Height of A}
\PY{+w}{  }\PY{n}{a2}\PY{+w}{ }\PY{o}{=}\PY{+w}{ }\PY{n}{atoi}\PY{p}{(}\PY{n}{argv}\PY{p}{[}\PY{l+m+mi}{2}\PY{p}{]}\PY{p}{)}\PY{p}{;}\PY{+w}{ }\PY{c+c1}{// Width of A}
\PY{+w}{  }\PY{n}{b1}\PY{+w}{ }\PY{o}{=}\PY{+w}{ }\PY{n}{a2}\PY{p}{;}\PY{+w}{            }\PY{c+c1}{// Height of B (must equal Width of A for multiplication)}
\PY{+w}{  }\PY{n}{b2}\PY{+w}{ }\PY{o}{=}\PY{+w}{ }\PY{n}{atoi}\PY{p}{(}\PY{n}{argv}\PY{p}{[}\PY{l+m+mi}{3}\PY{p}{]}\PY{p}{)}\PY{p}{;}\PY{+w}{ }\PY{c+c1}{// Width of B}

\PY{+w}{  }\PY{c+c1}{// Set dimensions for Matrix A and allocate host memory}
\PY{+w}{  }\PY{n}{A}\PY{p}{.}\PY{n}{height}\PY{+w}{ }\PY{o}{=}\PY{+w}{ }\PY{n}{a1}\PY{p}{;}
\PY{+w}{  }\PY{n}{A}\PY{p}{.}\PY{n}{width}\PY{+w}{ }\PY{o}{=}\PY{+w}{ }\PY{n}{a2}\PY{p}{;}
\PY{+w}{  }\PY{n}{A}\PY{p}{.}\PY{n}{elements}\PY{+w}{ }\PY{o}{=}\PY{+w}{ }\PY{p}{(}\PY{k+kt}{float}\PY{o}{*}\PY{p}{)}\PY{n}{malloc}\PY{p}{(}\PY{n}{A}\PY{p}{.}\PY{n}{width}\PY{+w}{ }\PY{o}{*}\PY{+w}{ }\PY{n}{A}\PY{p}{.}\PY{n}{height}\PY{+w}{ }\PY{o}{*}\PY{+w}{ }\PY{k}{sizeof}\PY{p}{(}\PY{k+kt}{float}\PY{p}{)}\PY{p}{)}\PY{p}{;}

\PY{+w}{  }\PY{c+c1}{// Set dimensions for Matrix B and allocate host memory}
\PY{+w}{  }\PY{n}{B}\PY{p}{.}\PY{n}{height}\PY{+w}{ }\PY{o}{=}\PY{+w}{ }\PY{n}{b1}\PY{p}{;}
\PY{+w}{  }\PY{n}{B}\PY{p}{.}\PY{n}{width}\PY{+w}{ }\PY{o}{=}\PY{+w}{ }\PY{n}{b2}\PY{p}{;}
\PY{+w}{  }\PY{n}{B}\PY{p}{.}\PY{n}{elements}\PY{+w}{ }\PY{o}{=}\PY{+w}{ }\PY{p}{(}\PY{k+kt}{float}\PY{o}{*}\PY{p}{)}\PY{n}{malloc}\PY{p}{(}\PY{n}{B}\PY{p}{.}\PY{n}{width}\PY{+w}{ }\PY{o}{*}\PY{+w}{ }\PY{n}{B}\PY{p}{.}\PY{n}{height}\PY{+w}{ }\PY{o}{*}\PY{+w}{ }\PY{k}{sizeof}\PY{p}{(}\PY{k+kt}{float}\PY{p}{)}\PY{p}{)}\PY{p}{;}

\PY{+w}{  }\PY{c+c1}{// Set dimensions for Matrix C (Output) and allocate host memory}
\PY{+w}{  }\PY{n}{C}\PY{p}{.}\PY{n}{height}\PY{+w}{ }\PY{o}{=}\PY{+w}{ }\PY{n}{A}\PY{p}{.}\PY{n}{height}\PY{p}{;}
\PY{+w}{  }\PY{n}{C}\PY{p}{.}\PY{n}{width}\PY{+w}{ }\PY{o}{=}\PY{+w}{ }\PY{n}{B}\PY{p}{.}\PY{n}{width}\PY{p}{;}
\PY{+w}{  }\PY{n}{C}\PY{p}{.}\PY{n}{elements}\PY{+w}{ }\PY{o}{=}\PY{+w}{ }\PY{p}{(}\PY{k+kt}{float}\PY{o}{*}\PY{p}{)}\PY{n}{malloc}\PY{p}{(}\PY{n}{C}\PY{p}{.}\PY{n}{width}\PY{+w}{ }\PY{o}{*}\PY{+w}{ }\PY{n}{C}\PY{p}{.}\PY{n}{height}\PY{+w}{ }\PY{o}{*}\PY{+w}{ }\PY{k}{sizeof}\PY{p}{(}\PY{k+kt}{float}\PY{p}{)}\PY{p}{)}\PY{p}{;}

\PY{+w}{  }\PY{c+c1}{// Populate Matrix A with random floats (0.0 or 1.0)}
\PY{+w}{  }\PY{k}{for}\PY{p}{(}\PY{k+kt}{int}\PY{+w}{ }\PY{n}{i}\PY{+w}{ }\PY{o}{=}\PY{+w}{ }\PY{l+m+mi}{0}\PY{p}{;}\PY{+w}{ }\PY{n}{i}\PY{+w}{ }\PY{o}{\PYZlt{}}\PY{+w}{ }\PY{n}{A}\PY{p}{.}\PY{n}{height}\PY{p}{;}\PY{+w}{ }\PY{n}{i}\PY{o}{+}\PY{o}{+}\PY{p}{)}
\PY{+w}{    }\PY{k}{for}\PY{p}{(}\PY{k+kt}{int}\PY{+w}{ }\PY{n}{j}\PY{+w}{ }\PY{o}{=}\PY{+w}{ }\PY{l+m+mi}{0}\PY{p}{;}\PY{+w}{ }\PY{n}{j}\PY{+w}{ }\PY{o}{\PYZlt{}}\PY{+w}{ }\PY{n}{A}\PY{p}{.}\PY{n}{width}\PY{p}{;}\PY{+w}{ }\PY{n}{j}\PY{o}{+}\PY{o}{+}\PY{p}{)}
\PY{+w}{      }\PY{n}{A}\PY{p}{.}\PY{n}{elements}\PY{p}{[}\PY{n}{i}\PY{o}{*}\PY{n}{A}\PY{p}{.}\PY{n}{width}\PY{+w}{ }\PY{o}{+}\PY{+w}{ }\PY{n}{j}\PY{p}{]}\PY{+w}{ }\PY{o}{=}\PY{+w}{ }\PY{p}{(}\PY{k+kt}{float}\PY{p}{)}\PY{p}{(}\PY{n}{rand}\PY{p}{(}\PY{p}{)}\PY{+w}{ }\PY{o}{\PYZpc{}}\PY{+w}{ }\PY{l+m+mi}{2}\PY{p}{)}\PY{p}{;}

\PY{+w}{  }\PY{c+c1}{// Populate Matrix B with random floats (0.0, 1.0, or 2.0)}
\PY{+w}{  }\PY{k}{for}\PY{p}{(}\PY{k+kt}{int}\PY{+w}{ }\PY{n}{i}\PY{+w}{ }\PY{o}{=}\PY{+w}{ }\PY{l+m+mi}{0}\PY{p}{;}\PY{+w}{ }\PY{n}{i}\PY{+w}{ }\PY{o}{\PYZlt{}}\PY{+w}{ }\PY{n}{B}\PY{p}{.}\PY{n}{height}\PY{p}{;}\PY{+w}{ }\PY{n}{i}\PY{o}{+}\PY{o}{+}\PY{p}{)}
\PY{+w}{    }\PY{k}{for}\PY{p}{(}\PY{k+kt}{int}\PY{+w}{ }\PY{n}{j}\PY{+w}{ }\PY{o}{=}\PY{+w}{ }\PY{l+m+mi}{0}\PY{p}{;}\PY{+w}{ }\PY{n}{j}\PY{+w}{ }\PY{o}{\PYZlt{}}\PY{+w}{ }\PY{n}{B}\PY{p}{.}\PY{n}{width}\PY{p}{;}\PY{+w}{ }\PY{n}{j}\PY{o}{+}\PY{o}{+}\PY{p}{)}
\PY{+w}{      }\PY{n}{B}\PY{p}{.}\PY{n}{elements}\PY{p}{[}\PY{n}{i}\PY{o}{*}\PY{n}{B}\PY{p}{.}\PY{n}{width}\PY{+w}{ }\PY{o}{+}\PY{+w}{ }\PY{n}{j}\PY{p}{]}\PY{+w}{ }\PY{o}{=}\PY{+w}{ }\PY{p}{(}\PY{k+kt}{float}\PY{p}{)}\PY{p}{(}\PY{n}{rand}\PY{p}{(}\PY{p}{)}\PY{+w}{ }\PY{o}{\PYZpc{}}\PY{+w}{ }\PY{l+m+mi}{3}\PY{p}{)}\PY{p}{;}

\PY{+w}{  }\PY{c+c1}{// Call the host function to execute the multiplication on the GPU}
\PY{+w}{  }\PY{n}{MatMul}\PY{p}{(}\PY{n}{A}\PY{p}{,}\PY{+w}{ }\PY{n}{B}\PY{p}{,}\PY{+w}{ }\PY{n}{C}\PY{p}{)}\PY{p}{;}

\PY{+w}{  }\PY{k}{return}\PY{+w}{ }\PY{l+m+mi}{0}\PY{p}{;}
\PY{p}{\PYZcb{}}

[Output of > MatrixMultiplication_global 1024 1024 1024]
CUDA malloc A: no error
Copy A to device: no error
CUDA malloc B: no error
Copy B to device: no error
CUDA malloc C: no error
Run kernel: no error
Time: 36.52512 ms
Copy C off of device: no error
\end{Verbatim}

\hrulefill

File: \underline{MatrixMultiplication\_local.cu}

\begin{Verbatim}[label=\makebox{\url{https://github.com/lucabaldini/cmepda/tree/master/slides/latex/snippets/MatrixMultiplication_shared.cu}},commandchars=\\\{\}]
\PY{c+cp}{\#}\PY{c+cp}{include}\PY{+w}{ }\PY{c+cpf}{\PYZlt{}stdlib.h\PYZgt{}}
\PY{c+cp}{\#}\PY{c+cp}{include}\PY{+w}{ }\PY{c+cpf}{\PYZlt{}stdio.h\PYZgt{}}
\PY{c+cp}{\#}\PY{c+cp}{include}\PY{+w}{ }\PY{c+cpf}{\PYZdq{}MatrixMultiplication\_shared.h\PYZdq{}}

\PY{c+c1}{// Matrix multiplication \PYZhy{} Host code (Shared Memory Version)}
\PY{c+c1}{// Matrix dimensions are assumed to be multiples of BLOCK\_SIZE}
\PY{k+kt}{void}\PY{+w}{ }\PY{n+nf}{MatMul}\PY{p}{(}\PY{k}{const}\PY{+w}{ }\PY{n}{Matrix}\PY{+w}{ }\PY{n}{A}\PY{p}{,}\PY{+w}{ }\PY{k}{const}\PY{+w}{ }\PY{n}{Matrix}\PY{+w}{ }\PY{n}{B}\PY{p}{,}\PY{+w}{ }\PY{n}{Matrix}\PY{+w}{ }\PY{n}{C}\PY{p}{)}\PY{+w}{ }\PY{p}{\PYZob{}}

\PY{+w}{  }\PY{c+c1}{// \PYZhy{}\PYZhy{}\PYZhy{} 1. ALLOCATE AND COPY MATRIX A TO DEVICE \PYZhy{}\PYZhy{}\PYZhy{}}
\PY{+w}{  }\PY{n}{Matrix}\PY{+w}{ }\PY{n}{d\_A}\PY{p}{;}
\PY{+w}{  }\PY{n}{d\_A}\PY{p}{.}\PY{n}{width}\PY{+w}{ }\PY{o}{=}\PY{+w}{ }\PY{n}{d\_A}\PY{p}{.}\PY{n}{stride}\PY{+w}{ }\PY{o}{=}\PY{+w}{ }\PY{n}{A}\PY{p}{.}\PY{n}{width}\PY{p}{;}
\PY{+w}{  }\PY{n}{d\_A}\PY{p}{.}\PY{n}{height}\PY{+w}{ }\PY{o}{=}\PY{+w}{ }\PY{n}{A}\PY{p}{.}\PY{n}{height}\PY{p}{;}

\PY{+w}{  }\PY{c+c1}{// Calculate the total size in bytes for Matrix A}
\PY{+w}{  }\PY{k+kt}{size\_t}\PY{+w}{ }\PY{n}{size}\PY{+w}{ }\PY{o}{=}\PY{+w}{ }\PY{n}{A}\PY{p}{.}\PY{n}{width}\PY{+w}{ }\PY{o}{*}\PY{+w}{ }\PY{n}{A}\PY{p}{.}\PY{n}{height}\PY{+w}{ }\PY{o}{*}\PY{+w}{ }\PY{k}{sizeof}\PY{p}{(}\PY{k+kt}{float}\PY{p}{)}\PY{p}{;}

\PY{+w}{  }\PY{c+c1}{// Allocate memory on the GPU (Device)}
\PY{+w}{  }\PY{n}{cudaError\_t}\PY{+w}{ }\PY{n}{err}\PY{+w}{ }\PY{o}{=}\PY{+w}{ }\PY{n}{cudaMalloc}\PY{p}{(}\PY{o}{\PYZam{}}\PY{n}{d\_A}\PY{p}{.}\PY{n}{elements}\PY{p}{,}\PY{+w}{ }\PY{n}{size}\PY{p}{)}\PY{p}{;}
\PY{+w}{  }\PY{n}{printf}\PY{p}{(}\PY{l+s}{\PYZdq{}}\PY{l+s}{CUDA malloc A: \PYZpc{}s}\PY{l+s+se}{\PYZbs{}n}\PY{l+s}{\PYZdq{}}\PY{p}{,}\PY{+w}{ }\PY{n}{cudaGetErrorString}\PY{p}{(}\PY{n}{err}\PY{p}{)}\PY{p}{)}\PY{p}{;}

\PY{+w}{  }\PY{c+c1}{// Copy data from Host (CPU) to Device (GPU)}
\PY{+w}{  }\PY{n}{cudaMemcpy}\PY{p}{(}\PY{n}{d\_A}\PY{p}{.}\PY{n}{elements}\PY{p}{,}\PY{+w}{ }\PY{n}{A}\PY{p}{.}\PY{n}{elements}\PY{p}{,}\PY{+w}{ }\PY{n}{size}\PY{p}{,}\PY{+w}{ }\PY{n}{cudaMemcpyHostToDevice}\PY{p}{)}\PY{p}{;}

\PY{+w}{  }\PY{c+c1}{// \PYZhy{}\PYZhy{}\PYZhy{} 2. ALLOCATE AND COPY MATRIX B TO DEVICE \PYZhy{}\PYZhy{}\PYZhy{}}
\PY{+w}{  }\PY{n}{Matrix}\PY{+w}{ }\PY{n}{d\_B}\PY{p}{;}
\PY{+w}{  }\PY{n}{d\_B}\PY{p}{.}\PY{n}{width}\PY{+w}{ }\PY{o}{=}\PY{+w}{ }\PY{n}{d\_B}\PY{p}{.}\PY{n}{stride}\PY{+w}{ }\PY{o}{=}\PY{+w}{ }\PY{n}{B}\PY{p}{.}\PY{n}{width}\PY{p}{;}
\PY{+w}{  }\PY{n}{d\_B}\PY{p}{.}\PY{n}{height}\PY{+w}{ }\PY{o}{=}\PY{+w}{ }\PY{n}{B}\PY{p}{.}\PY{n}{height}\PY{p}{;}
\PY{+w}{  }\PY{n}{size}\PY{+w}{ }\PY{o}{=}\PY{+w}{ }\PY{n}{B}\PY{p}{.}\PY{n}{width}\PY{+w}{ }\PY{o}{*}\PY{+w}{ }\PY{n}{B}\PY{p}{.}\PY{n}{height}\PY{+w}{ }\PY{o}{*}\PY{+w}{ }\PY{k}{sizeof}\PY{p}{(}\PY{k+kt}{float}\PY{p}{)}\PY{p}{;}

\PY{+w}{  }\PY{n}{err}\PY{+w}{ }\PY{o}{=}\PY{+w}{ }\PY{n}{cudaMalloc}\PY{p}{(}\PY{o}{\PYZam{}}\PY{n}{d\_B}\PY{p}{.}\PY{n}{elements}\PY{p}{,}\PY{+w}{ }\PY{n}{size}\PY{p}{)}\PY{p}{;}
\PY{+w}{  }\PY{n}{printf}\PY{p}{(}\PY{l+s}{\PYZdq{}}\PY{l+s}{CUDA malloc B: \PYZpc{}s}\PY{l+s+se}{\PYZbs{}n}\PY{l+s}{\PYZdq{}}\PY{p}{,}\PY{+w}{ }\PY{n}{cudaGetErrorString}\PY{p}{(}\PY{n}{err}\PY{p}{)}\PY{p}{)}\PY{p}{;}
\PY{+w}{  }\PY{n}{cudaMemcpy}\PY{p}{(}\PY{n}{d\_B}\PY{p}{.}\PY{n}{elements}\PY{p}{,}\PY{+w}{ }\PY{n}{B}\PY{p}{.}\PY{n}{elements}\PY{p}{,}\PY{+w}{ }\PY{n}{size}\PY{p}{,}\PY{+w}{ }\PY{n}{cudaMemcpyHostToDevice}\PY{p}{)}\PY{p}{;}

\PY{+w}{  }\PY{c+c1}{// \PYZhy{}\PYZhy{}\PYZhy{} 3. ALLOCATE MATRIX C ON DEVICE \PYZhy{}\PYZhy{}\PYZhy{}}
\PY{+w}{  }\PY{n}{Matrix}\PY{+w}{ }\PY{n}{d\_C}\PY{p}{;}
\PY{+w}{  }\PY{n}{d\_C}\PY{p}{.}\PY{n}{width}\PY{+w}{ }\PY{o}{=}\PY{+w}{ }\PY{n}{d\_C}\PY{p}{.}\PY{n}{stride}\PY{+w}{ }\PY{o}{=}\PY{+w}{ }\PY{n}{C}\PY{p}{.}\PY{n}{width}\PY{p}{;}
\PY{+w}{  }\PY{n}{d\_C}\PY{p}{.}\PY{n}{height}\PY{+w}{ }\PY{o}{=}\PY{+w}{ }\PY{n}{C}\PY{p}{.}\PY{n}{height}\PY{p}{;}
\PY{+w}{  }\PY{n}{size}\PY{+w}{ }\PY{o}{=}\PY{+w}{ }\PY{n}{C}\PY{p}{.}\PY{n}{width}\PY{+w}{ }\PY{o}{*}\PY{+w}{ }\PY{n}{C}\PY{p}{.}\PY{n}{height}\PY{+w}{ }\PY{o}{*}\PY{+w}{ }\PY{k}{sizeof}\PY{p}{(}\PY{k+kt}{float}\PY{p}{)}\PY{p}{;}

\PY{+w}{  }\PY{n}{err}\PY{+w}{ }\PY{o}{=}\PY{+w}{ }\PY{n}{cudaMalloc}\PY{p}{(}\PY{o}{\PYZam{}}\PY{n}{d\_C}\PY{p}{.}\PY{n}{elements}\PY{p}{,}\PY{+w}{ }\PY{n}{size}\PY{p}{)}\PY{p}{;}
\PY{+w}{  }\PY{n}{printf}\PY{p}{(}\PY{l+s}{\PYZdq{}}\PY{l+s}{CUDA malloc C: \PYZpc{}s}\PY{l+s+se}{\PYZbs{}n}\PY{l+s}{\PYZdq{}}\PY{p}{,}\PY{+w}{ }\PY{n}{cudaGetErrorString}\PY{p}{(}\PY{n}{err}\PY{p}{)}\PY{p}{)}\PY{p}{;}

\PY{+w}{  }\PY{c+c1}{// \PYZhy{}\PYZhy{}\PYZhy{} 4. TIMING SETUP \PYZhy{}\PYZhy{}\PYZhy{}}
\PY{+w}{  }\PY{k+kt}{float}\PY{+w}{ }\PY{n}{time}\PY{p}{;}
\PY{+w}{  }\PY{n}{cudaEvent\_t}\PY{+w}{ }\PY{n}{start}\PY{p}{,}\PY{+w}{ }\PY{n}{stop}\PY{p}{;}
\PY{+w}{  }\PY{n}{cudaEventCreate}\PY{p}{(}\PY{o}{\PYZam{}}\PY{n}{start}\PY{p}{)}\PY{p}{;}
\PY{+w}{  }\PY{n}{cudaEventCreate}\PY{p}{(}\PY{o}{\PYZam{}}\PY{n}{stop}\PY{p}{)}\PY{p}{;}

\PY{+w}{  }\PY{c+c1}{// Start the timer}
\PY{+w}{  }\PY{n}{cudaEventRecord}\PY{p}{(}\PY{n}{start}\PY{p}{)}\PY{p}{;}

\PY{+w}{  }\PY{c+c1}{// \PYZhy{}\PYZhy{}\PYZhy{} 5. GRID AND BLOCK CONFIGURATION \PYZhy{}\PYZhy{}\PYZhy{}}
\PY{+w}{  }\PY{c+c1}{// Define the number of threads per block in 2D (BLOCK\_SIZE x BLOCK\_SIZE)}
\PY{+w}{  }\PY{k+kt}{dim3}\PY{+w}{ }\PY{n}{dimBlock}\PY{p}{(}\PY{n}{BLOCK\_SIZE}\PY{p}{,}\PY{+w}{ }\PY{n}{BLOCK\_SIZE}\PY{p}{)}\PY{p}{;}

\PY{+w}{  }\PY{c+c1}{// Calculate the number of blocks needed in the grid}
\PY{+w}{  }\PY{k+kt}{dim3}\PY{+w}{ }\PY{n}{dimGrid}\PY{p}{(}\PY{n}{B}\PY{p}{.}\PY{n}{width}\PY{+w}{ }\PY{o}{/}\PY{+w}{ }\PY{n}{dimBlock}\PY{p}{.}\PY{n}{x}\PY{p}{,}\PY{+w}{ }\PY{n}{A}\PY{p}{.}\PY{n}{height}\PY{+w}{ }\PY{o}{/}\PY{+w}{ }\PY{n}{dimBlock}\PY{p}{.}\PY{n}{y}\PY{p}{)}\PY{p}{;}

\PY{+w}{  }\PY{c+c1}{// Launch the kernel}
\PY{+w}{  }\PY{n}{MatMulKernel}\PY{o}{\PYZlt{}}\PY{o}{\PYZlt{}}\PY{o}{\PYZlt{}}\PY{n}{dimGrid}\PY{p}{,}\PY{+w}{ }\PY{n}{dimBlock}\PY{o}{\PYZgt{}}\PY{o}{\PYZgt{}}\PY{o}{\PYZgt{}}\PY{p}{(}\PY{n}{d\_A}\PY{p}{,}\PY{+w}{ }\PY{n}{d\_B}\PY{p}{,}\PY{+w}{ }\PY{n}{d\_C}\PY{p}{)}\PY{p}{;}

\PY{+w}{  }\PY{c+c1}{// Wait for all threads to finish}
\PY{+w}{  }\PY{n}{err}\PY{+w}{ }\PY{o}{=}\PY{+w}{ }\PY{n}{cudaDeviceSynchronize}\PY{p}{(}\PY{p}{)}\PY{p}{;}

\PY{+w}{  }\PY{c+c1}{// \PYZhy{}\PYZhy{}\PYZhy{} 6. STOP TIMER AND PRINT \PYZhy{}\PYZhy{}\PYZhy{}}
\PY{+w}{  }\PY{n}{cudaEventRecord}\PY{p}{(}\PY{n}{stop}\PY{p}{)}\PY{p}{;}
\PY{+w}{  }\PY{n}{cudaEventSynchronize}\PY{p}{(}\PY{n}{stop}\PY{p}{)}\PY{p}{;}
\PY{+w}{  }\PY{n}{cudaEventElapsedTime}\PY{p}{(}\PY{o}{\PYZam{}}\PY{n}{time}\PY{p}{,}\PY{+w}{ }\PY{n}{start}\PY{p}{,}\PY{+w}{ }\PY{n}{stop}\PY{p}{)}\PY{p}{;}

\PY{+w}{  }\PY{n}{printf}\PY{p}{(}\PY{l+s}{\PYZdq{}}\PY{l+s}{Run kernel: \PYZpc{}s}\PY{l+s+se}{\PYZbs{}n}\PY{l+s}{\PYZdq{}}\PY{p}{,}\PY{+w}{ }\PY{n}{cudaGetErrorString}\PY{p}{(}\PY{n}{err}\PY{p}{)}\PY{p}{)}\PY{p}{;}
\PY{+w}{  }\PY{n}{printf}\PY{p}{(}\PY{l+s}{\PYZdq{}}\PY{l+s}{Time: \PYZpc{}3.5f ms}\PY{l+s+se}{\PYZbs{}n}\PY{l+s}{\PYZdq{}}\PY{p}{,}\PY{+w}{ }\PY{n}{time}\PY{p}{)}\PY{p}{;}

\PY{+w}{  }\PY{c+c1}{// \PYZhy{}\PYZhy{}\PYZhy{} 7. RETRIEVE RESULTS AND CLEANUP \PYZhy{}\PYZhy{}\PYZhy{}}
\PY{+w}{  }\PY{c+c1}{// Copy the computed matrix C back to the Host}
\PY{+w}{  }\PY{n}{err}\PY{+w}{ }\PY{o}{=}\PY{+w}{ }\PY{n}{cudaMemcpy}\PY{p}{(}\PY{n}{C}\PY{p}{.}\PY{n}{elements}\PY{p}{,}\PY{+w}{ }\PY{n}{d\_C}\PY{p}{.}\PY{n}{elements}\PY{p}{,}\PY{+w}{ }\PY{n}{size}\PY{p}{,}\PY{+w}{ }\PY{n}{cudaMemcpyDeviceToHost}\PY{p}{)}\PY{p}{;}
\PY{+w}{  }\PY{n}{printf}\PY{p}{(}\PY{l+s}{\PYZdq{}}\PY{l+s}{Copy C off of device: \PYZpc{}s}\PY{l+s+se}{\PYZbs{}n}\PY{l+s}{\PYZdq{}}\PY{p}{,}\PY{+w}{ }\PY{n}{cudaGetErrorString}\PY{p}{(}\PY{n}{err}\PY{p}{)}\PY{p}{)}\PY{p}{;}

\PY{+w}{  }\PY{c+c1}{// Free the device memory to prevent memory leaks}
\PY{+w}{  }\PY{n}{cudaFree}\PY{p}{(}\PY{n}{d\_A}\PY{p}{.}\PY{n}{elements}\PY{p}{)}\PY{p}{;}
\PY{+w}{  }\PY{n}{cudaFree}\PY{p}{(}\PY{n}{d\_B}\PY{p}{.}\PY{n}{elements}\PY{p}{)}\PY{p}{;}
\PY{+w}{  }\PY{n}{cudaFree}\PY{p}{(}\PY{n}{d\_C}\PY{p}{.}\PY{n}{elements}\PY{p}{)}\PY{p}{;}
\PY{p}{\PYZcb{}}

\PY{c+c1}{// ===============================================================}
\PY{c+c1}{// \PYZhy{}\PYZhy{}\PYZhy{} DEVICE HELPER FUNCTIONS \PYZhy{}\PYZhy{}\PYZhy{}}
\PY{c+c1}{// ===============================================================}

\PY{c+c1}{// Get a matrix element}
\PY{k+kt}{\_\_device\_\_}\PY{+w}{ }\PY{k+kt}{float}\PY{+w}{ }\PY{n}{GetElement}\PY{p}{(}\PY{k}{const}\PY{+w}{ }\PY{n}{Matrix}\PY{+w}{ }\PY{n}{A}\PY{p}{,}\PY{+w}{ }\PY{k+kt}{int}\PY{+w}{ }\PY{n}{row}\PY{p}{,}\PY{+w}{ }\PY{k+kt}{int}\PY{+w}{ }\PY{n}{col}\PY{p}{)}\PY{+w}{ }\PY{p}{\PYZob{}}
\PY{+w}{  }\PY{k}{return}\PY{+w}{ }\PY{n}{A}\PY{p}{.}\PY{n}{elements}\PY{p}{[}\PY{n}{row}\PY{+w}{ }\PY{o}{*}\PY{+w}{ }\PY{n}{A}\PY{p}{.}\PY{n}{stride}\PY{+w}{ }\PY{o}{+}\PY{+w}{ }\PY{n}{col}\PY{p}{]}\PY{p}{;}
\PY{p}{\PYZcb{}}

\PY{c+c1}{// Set a matrix element}
\PY{k+kt}{\_\_device\_\_}\PY{+w}{ }\PY{k+kt}{void}\PY{+w}{ }\PY{n}{SetElement}\PY{p}{(}\PY{n}{Matrix}\PY{+w}{ }\PY{n}{A}\PY{p}{,}\PY{+w}{ }\PY{k+kt}{int}\PY{+w}{ }\PY{n}{row}\PY{p}{,}\PY{+w}{ }\PY{k+kt}{int}\PY{+w}{ }\PY{n}{col}\PY{p}{,}\PY{+w}{ }\PY{k+kt}{float}\PY{+w}{ }\PY{n}{value}\PY{p}{)}\PY{+w}{ }\PY{p}{\PYZob{}}
\PY{+w}{  }\PY{n}{A}\PY{p}{.}\PY{n}{elements}\PY{p}{[}\PY{n}{row}\PY{+w}{ }\PY{o}{*}\PY{+w}{ }\PY{n}{A}\PY{p}{.}\PY{n}{stride}\PY{+w}{ }\PY{o}{+}\PY{+w}{ }\PY{n}{col}\PY{p}{]}\PY{+w}{ }\PY{o}{=}\PY{+w}{ }\PY{n}{value}\PY{p}{;}
\PY{p}{\PYZcb{}}

\PY{c+c1}{// Get the BLOCK\_SIZE x BLOCK\_SIZE sub\PYZhy{}matrix Asub of A that is}
\PY{c+c1}{// located col sub\PYZhy{}matrices to the right and row sub\PYZhy{}matrices down}
\PY{c+c1}{// from the upper\PYZhy{}left corner of A}
\PY{k+kt}{\_\_device\_\_}\PY{+w}{ }\PY{n}{Matrix}\PY{+w}{ }\PY{n}{GetSubMatrix}\PY{p}{(}\PY{n}{Matrix}\PY{+w}{ }\PY{n}{A}\PY{p}{,}\PY{+w}{ }\PY{k+kt}{int}\PY{+w}{ }\PY{n}{row}\PY{p}{,}\PY{+w}{ }\PY{k+kt}{int}\PY{+w}{ }\PY{n}{col}\PY{p}{)}\PY{+w}{ }\PY{p}{\PYZob{}}
\PY{+w}{  }\PY{n}{Matrix}\PY{+w}{ }\PY{n}{Asub}\PY{p}{;}
\PY{+w}{  }\PY{n}{Asub}\PY{p}{.}\PY{n}{width}\PY{+w}{ }\PY{o}{=}\PY{+w}{ }\PY{n}{BLOCK\_SIZE}\PY{p}{;}
\PY{+w}{  }\PY{n}{Asub}\PY{p}{.}\PY{n}{height}\PY{+w}{ }\PY{o}{=}\PY{+w}{ }\PY{n}{BLOCK\_SIZE}\PY{p}{;}
\PY{+w}{  }\PY{n}{Asub}\PY{p}{.}\PY{n}{stride}\PY{+w}{ }\PY{o}{=}\PY{+w}{ }\PY{n}{A}\PY{p}{.}\PY{n}{stride}\PY{p}{;}
\PY{+w}{  }\PY{n}{Asub}\PY{p}{.}\PY{n}{elements}\PY{+w}{ }\PY{o}{=}\PY{+w}{ }\PY{o}{\PYZam{}}\PY{n}{A}\PY{p}{.}\PY{n}{elements}\PY{p}{[}\PY{n}{A}\PY{p}{.}\PY{n}{stride}\PY{+w}{ }\PY{o}{*}\PY{+w}{ }\PY{n}{BLOCK\_SIZE}\PY{+w}{ }\PY{o}{*}\PY{+w}{ }\PY{n}{row}\PY{+w}{ }\PY{o}{+}\PY{+w}{ }\PY{n}{BLOCK\_SIZE}\PY{+w}{ }\PY{o}{*}\PY{+w}{ }\PY{n}{col}\PY{p}{]}\PY{p}{;}
\PY{+w}{  }\PY{k}{return}\PY{+w}{ }\PY{n}{Asub}\PY{p}{;}
\PY{p}{\PYZcb{}}

\PY{c+c1}{// ===============================================================}
\PY{c+c1}{// \PYZhy{}\PYZhy{}\PYZhy{} KERNEL FUNCTION \PYZhy{}\PYZhy{}\PYZhy{}}
\PY{c+c1}{// ===============================================================}

\PY{c+c1}{// Matrix multiplication kernel called by MatMul()}
\PY{k+kr}{\_\_global\_\_}\PY{+w}{ }\PY{k+kt}{void}\PY{+w}{ }\PY{n}{MatMulKernel}\PY{p}{(}\PY{n}{Matrix}\PY{+w}{ }\PY{n}{A}\PY{p}{,}\PY{+w}{ }\PY{n}{Matrix}\PY{+w}{ }\PY{n}{B}\PY{p}{,}\PY{+w}{ }\PY{n}{Matrix}\PY{+w}{ }\PY{n}{C}\PY{p}{)}\PY{+w}{ }\PY{p}{\PYZob{}}

\PY{+w}{  }\PY{c+c1}{// Calculate the global block row and column index}
\PY{+w}{  }\PY{c+c1}{// blockIdx: Which block we are in}
\PY{+w}{  }\PY{k+kt}{int}\PY{+w}{ }\PY{n}{blockRow}\PY{+w}{ }\PY{o}{=}\PY{+w}{ }\PY{n+nb}{blockIdx}\PY{p}{.}\PY{n}{y}\PY{p}{;}
\PY{+w}{  }\PY{k+kt}{int}\PY{+w}{ }\PY{n}{blockCol}\PY{+w}{ }\PY{o}{=}\PY{+w}{ }\PY{n+nb}{blockIdx}\PY{p}{.}\PY{n}{x}\PY{p}{;}

\PY{+w}{  }\PY{c+c1}{// Each thread block computes one sub\PYZhy{}matrix Csub of C}
\PY{+w}{  }\PY{n}{Matrix}\PY{+w}{ }\PY{n}{Csub}\PY{+w}{ }\PY{o}{=}\PY{+w}{ }\PY{n}{GetSubMatrix}\PY{p}{(}\PY{n}{C}\PY{p}{,}\PY{+w}{ }\PY{n}{blockRow}\PY{p}{,}\PY{+w}{ }\PY{n}{blockCol}\PY{p}{)}\PY{p}{;}

\PY{+w}{  }\PY{c+c1}{// Variable to accumulate the dot product sum for this specific thread\PYZsq{}s element}
\PY{+w}{  }\PY{k+kt}{float}\PY{+w}{ }\PY{n}{Cvalue}\PY{+w}{ }\PY{o}{=}\PY{+w}{ }\PY{l+m+mf}{0.0}\PY{p}{;}

\PY{+w}{  }\PY{c+c1}{// Thread row and column within Csub}
\PY{+w}{  }\PY{c+c1}{// threadIdx: Which thread we are inside the block}
\PY{+w}{  }\PY{k+kt}{int}\PY{+w}{ }\PY{n}{row}\PY{+w}{ }\PY{o}{=}\PY{+w}{ }\PY{n+nb}{threadIdx}\PY{p}{.}\PY{n}{y}\PY{p}{;}
\PY{+w}{  }\PY{k+kt}{int}\PY{+w}{ }\PY{n}{col}\PY{+w}{ }\PY{o}{=}\PY{+w}{ }\PY{n+nb}{threadIdx}\PY{p}{.}\PY{n}{x}\PY{p}{;}

\PY{+w}{  }\PY{c+c1}{// Loop over all the sub\PYZhy{}matrices of A and B that are}
\PY{+w}{  }\PY{c+c1}{// required to compute Csub}
\PY{+w}{  }\PY{k}{for}\PY{+w}{ }\PY{p}{(}\PY{k+kt}{int}\PY{+w}{ }\PY{n}{m}\PY{+w}{ }\PY{o}{=}\PY{+w}{ }\PY{l+m+mi}{0}\PY{p}{;}\PY{+w}{ }\PY{n}{m}\PY{+w}{ }\PY{o}{\PYZlt{}}\PY{+w}{ }\PY{p}{(}\PY{n}{A}\PY{p}{.}\PY{n}{width}\PY{+w}{ }\PY{o}{/}\PY{+w}{ }\PY{n}{BLOCK\_SIZE}\PY{p}{)}\PY{p}{;}\PY{+w}{ }\PY{o}{+}\PY{o}{+}\PY{n}{m}\PY{p}{)}\PY{+w}{ }\PY{p}{\PYZob{}}

\PY{+w}{    }\PY{c+c1}{// Get sub\PYZhy{}matrix Asub of A and Bsub of B}
\PY{+w}{    }\PY{n}{Matrix}\PY{+w}{ }\PY{n}{Asub}\PY{+w}{ }\PY{o}{=}\PY{+w}{ }\PY{n}{GetSubMatrix}\PY{p}{(}\PY{n}{A}\PY{p}{,}\PY{+w}{ }\PY{n}{blockRow}\PY{p}{,}\PY{+w}{ }\PY{n}{m}\PY{p}{)}\PY{p}{;}
\PY{+w}{    }\PY{n}{Matrix}\PY{+w}{ }\PY{n}{Bsub}\PY{+w}{ }\PY{o}{=}\PY{+w}{ }\PY{n}{GetSubMatrix}\PY{p}{(}\PY{n}{B}\PY{p}{,}\PY{+w}{ }\PY{n}{m}\PY{p}{,}\PY{+w}{ }\PY{n}{blockCol}\PY{p}{)}\PY{p}{;}

\PY{+w}{    }\PY{c+c1}{// Shared memory used to store Asub and Bsub respectively}
\PY{+w}{    }\PY{k+kt}{\_\_shared\_\_}\PY{+w}{ }\PY{k+kt}{float}\PY{+w}{ }\PY{n}{As}\PY{p}{[}\PY{n}{BLOCK\_SIZE}\PY{p}{]}\PY{p}{[}\PY{n}{BLOCK\_SIZE}\PY{p}{]}\PY{p}{;}
\PY{+w}{    }\PY{k+kt}{\_\_shared\_\_}\PY{+w}{ }\PY{k+kt}{float}\PY{+w}{ }\PY{n}{Bs}\PY{p}{[}\PY{n}{BLOCK\_SIZE}\PY{p}{]}\PY{p}{[}\PY{n}{BLOCK\_SIZE}\PY{p}{]}\PY{p}{;}

\PY{+w}{    }\PY{c+c1}{// Load Asub and Bsub from device memory to shared memory}
\PY{+w}{    }\PY{c+c1}{// Each thread loads one element of each sub\PYZhy{}matrix}
\PY{+w}{    }\PY{n}{As}\PY{p}{[}\PY{n}{row}\PY{p}{]}\PY{p}{[}\PY{n}{col}\PY{p}{]}\PY{+w}{ }\PY{o}{=}\PY{+w}{ }\PY{n}{GetElement}\PY{p}{(}\PY{n}{Asub}\PY{p}{,}\PY{+w}{ }\PY{n}{row}\PY{p}{,}\PY{+w}{ }\PY{n}{col}\PY{p}{)}\PY{p}{;}
\PY{+w}{    }\PY{n}{Bs}\PY{p}{[}\PY{n}{row}\PY{p}{]}\PY{p}{[}\PY{n}{col}\PY{p}{]}\PY{+w}{ }\PY{o}{=}\PY{+w}{ }\PY{n}{GetElement}\PY{p}{(}\PY{n}{Bsub}\PY{p}{,}\PY{+w}{ }\PY{n}{row}\PY{p}{,}\PY{+w}{ }\PY{n}{col}\PY{p}{)}\PY{p}{;}

\PY{+w}{    }\PY{c+c1}{// Synchronize to make sure the sub\PYZhy{}matrices are fully loaded}
\PY{+w}{    }\PY{c+c1}{// before starting the computation}
\PY{+w}{    }\PY{n+nf}{\_\_syncthreads}\PY{p}{(}\PY{p}{)}\PY{p}{;}

\PY{+w}{    }\PY{c+c1}{// Multiply Asub and Bsub together and accumulate the results}
\PY{+w}{    }\PY{k}{for}\PY{+w}{ }\PY{p}{(}\PY{k+kt}{int}\PY{+w}{ }\PY{n}{e}\PY{+w}{ }\PY{o}{=}\PY{+w}{ }\PY{l+m+mi}{0}\PY{p}{;}\PY{+w}{ }\PY{n}{e}\PY{+w}{ }\PY{o}{\PYZlt{}}\PY{+w}{ }\PY{n}{BLOCK\_SIZE}\PY{p}{;}\PY{+w}{ }\PY{o}{+}\PY{o}{+}\PY{n}{e}\PY{p}{)}\PY{+w}{ }\PY{p}{\PYZob{}}
\PY{+w}{      }\PY{n}{Cvalue}\PY{+w}{ }\PY{o}{+}\PY{o}{=}\PY{+w}{ }\PY{n}{As}\PY{p}{[}\PY{n}{row}\PY{p}{]}\PY{p}{[}\PY{n}{e}\PY{p}{]}\PY{+w}{ }\PY{o}{*}\PY{+w}{ }\PY{n}{Bs}\PY{p}{[}\PY{n}{e}\PY{p}{]}\PY{p}{[}\PY{n}{col}\PY{p}{]}\PY{p}{;}
\PY{+w}{    }\PY{p}{\PYZcb{}}

\PY{+w}{    }\PY{c+c1}{// Synchronize to make sure that the preceding computation is done}
\PY{+w}{    }\PY{c+c1}{// before loading two new sub\PYZhy{}matrices of A and B in the next iteration}
\PY{+w}{    }\PY{n+nf}{\_\_syncthreads}\PY{p}{(}\PY{p}{)}\PY{p}{;}
\PY{+w}{  }\PY{p}{\PYZcb{}}

\PY{+w}{  }\PY{c+c1}{// Write the final accumulated value to the sub\PYZhy{}matrix Csub in device memory}
\PY{+w}{  }\PY{n}{SetElement}\PY{p}{(}\PY{n}{Csub}\PY{p}{,}\PY{+w}{ }\PY{n}{row}\PY{p}{,}\PY{+w}{ }\PY{n}{col}\PY{p}{,}\PY{+w}{ }\PY{n}{Cvalue}\PY{p}{)}\PY{p}{;}
\PY{p}{\PYZcb{}}

\PY{c+c1}{// Usage: .\PYZbs{}MatrixMultiplication\_shared a1 a2 b2}
\PY{k+kt}{int}\PY{+w}{ }\PY{n}{main}\PY{p}{(}\PY{k+kt}{int}\PY{+w}{ }\PY{n}{argc}\PY{p}{,}\PY{+w}{ }\PY{k+kt}{char}\PY{o}{*}\PY{+w}{ }\PY{n}{argv}\PY{p}{[}\PY{p}{]}\PY{p}{)}\PY{p}{\PYZob{}}
\PY{+w}{  }\PY{n}{Matrix}\PY{+w}{ }\PY{n}{A}\PY{p}{,}\PY{+w}{ }\PY{n}{B}\PY{p}{,}\PY{+w}{ }\PY{n}{C}\PY{p}{;}
\PY{+w}{  }\PY{k+kt}{int}\PY{+w}{ }\PY{n}{a1}\PY{p}{,}\PY{+w}{ }\PY{n}{a2}\PY{p}{,}\PY{+w}{ }\PY{n}{b1}\PY{p}{,}\PY{+w}{ }\PY{n}{b2}\PY{p}{;}

\PY{+w}{  }\PY{c+c1}{// Read matrix dimensions from command line arguments}
\PY{+w}{  }\PY{n}{a1}\PY{+w}{ }\PY{o}{=}\PY{+w}{ }\PY{n}{atoi}\PY{p}{(}\PY{n}{argv}\PY{p}{[}\PY{l+m+mi}{1}\PY{p}{]}\PY{p}{)}\PY{p}{;}\PY{+w}{ }\PY{c+c1}{// Height of A}
\PY{+w}{  }\PY{n}{a2}\PY{+w}{ }\PY{o}{=}\PY{+w}{ }\PY{n}{atoi}\PY{p}{(}\PY{n}{argv}\PY{p}{[}\PY{l+m+mi}{2}\PY{p}{]}\PY{p}{)}\PY{p}{;}\PY{+w}{ }\PY{c+c1}{// Width of A}
\PY{+w}{  }\PY{n}{b1}\PY{+w}{ }\PY{o}{=}\PY{+w}{ }\PY{n}{a2}\PY{p}{;}\PY{+w}{            }\PY{c+c1}{// Height of B (must equal Width of A for multiplication)}
\PY{+w}{  }\PY{n}{b2}\PY{+w}{ }\PY{o}{=}\PY{+w}{ }\PY{n}{atoi}\PY{p}{(}\PY{n}{argv}\PY{p}{[}\PY{l+m+mi}{3}\PY{p}{]}\PY{p}{)}\PY{p}{;}\PY{+w}{ }\PY{c+c1}{// Width of B}

\PY{+w}{  }\PY{c+c1}{// Set dimensions for Matrix A and allocate host memory}
\PY{+w}{  }\PY{n}{A}\PY{p}{.}\PY{n}{height}\PY{+w}{ }\PY{o}{=}\PY{+w}{ }\PY{n}{a1}\PY{p}{;}
\PY{+w}{  }\PY{n}{A}\PY{p}{.}\PY{n}{width}\PY{+w}{ }\PY{o}{=}\PY{+w}{ }\PY{n}{a2}\PY{p}{;}
\PY{+w}{  }\PY{n}{A}\PY{p}{.}\PY{n}{elements}\PY{+w}{ }\PY{o}{=}\PY{+w}{ }\PY{p}{(}\PY{k+kt}{float}\PY{o}{*}\PY{p}{)}\PY{n}{malloc}\PY{p}{(}\PY{n}{A}\PY{p}{.}\PY{n}{width}\PY{+w}{ }\PY{o}{*}\PY{+w}{ }\PY{n}{A}\PY{p}{.}\PY{n}{height}\PY{+w}{ }\PY{o}{*}\PY{+w}{ }\PY{k}{sizeof}\PY{p}{(}\PY{k+kt}{float}\PY{p}{)}\PY{p}{)}\PY{p}{;}

\PY{+w}{  }\PY{c+c1}{// Set dimensions for Matrix B and allocate host memory}
\PY{+w}{  }\PY{n}{B}\PY{p}{.}\PY{n}{height}\PY{+w}{ }\PY{o}{=}\PY{+w}{ }\PY{n}{b1}\PY{p}{;}
\PY{+w}{  }\PY{n}{B}\PY{p}{.}\PY{n}{width}\PY{+w}{ }\PY{o}{=}\PY{+w}{ }\PY{n}{b2}\PY{p}{;}
\PY{+w}{  }\PY{n}{B}\PY{p}{.}\PY{n}{elements}\PY{+w}{ }\PY{o}{=}\PY{+w}{ }\PY{p}{(}\PY{k+kt}{float}\PY{o}{*}\PY{p}{)}\PY{n}{malloc}\PY{p}{(}\PY{n}{B}\PY{p}{.}\PY{n}{width}\PY{+w}{ }\PY{o}{*}\PY{+w}{ }\PY{n}{B}\PY{p}{.}\PY{n}{height}\PY{+w}{ }\PY{o}{*}\PY{+w}{ }\PY{k}{sizeof}\PY{p}{(}\PY{k+kt}{float}\PY{p}{)}\PY{p}{)}\PY{p}{;}

\PY{+w}{  }\PY{c+c1}{// Set dimensions for Matrix C (Output) and allocate host memory}
\PY{+w}{  }\PY{n}{C}\PY{p}{.}\PY{n}{height}\PY{+w}{ }\PY{o}{=}\PY{+w}{ }\PY{n}{A}\PY{p}{.}\PY{n}{height}\PY{p}{;}
\PY{+w}{  }\PY{n}{C}\PY{p}{.}\PY{n}{width}\PY{+w}{ }\PY{o}{=}\PY{+w}{ }\PY{n}{B}\PY{p}{.}\PY{n}{width}\PY{p}{;}
\PY{+w}{  }\PY{n}{C}\PY{p}{.}\PY{n}{elements}\PY{+w}{ }\PY{o}{=}\PY{+w}{ }\PY{p}{(}\PY{k+kt}{float}\PY{o}{*}\PY{p}{)}\PY{n}{malloc}\PY{p}{(}\PY{n}{C}\PY{p}{.}\PY{n}{width}\PY{+w}{ }\PY{o}{*}\PY{+w}{ }\PY{n}{C}\PY{p}{.}\PY{n}{height}\PY{+w}{ }\PY{o}{*}\PY{+w}{ }\PY{k}{sizeof}\PY{p}{(}\PY{k+kt}{float}\PY{p}{)}\PY{p}{)}\PY{p}{;}

\PY{+w}{  }\PY{c+c1}{// Populate Matrix A with random floats (0.0, 1.0, or 2.0)}
\PY{+w}{  }\PY{k}{for}\PY{p}{(}\PY{k+kt}{int}\PY{+w}{ }\PY{n}{i}\PY{+w}{ }\PY{o}{=}\PY{+w}{ }\PY{l+m+mi}{0}\PY{p}{;}\PY{+w}{ }\PY{n}{i}\PY{+w}{ }\PY{o}{\PYZlt{}}\PY{+w}{ }\PY{n}{A}\PY{p}{.}\PY{n}{height}\PY{p}{;}\PY{+w}{ }\PY{n}{i}\PY{o}{+}\PY{o}{+}\PY{p}{)}
\PY{+w}{    }\PY{k}{for}\PY{p}{(}\PY{k+kt}{int}\PY{+w}{ }\PY{n}{j}\PY{+w}{ }\PY{o}{=}\PY{+w}{ }\PY{l+m+mi}{0}\PY{p}{;}\PY{+w}{ }\PY{n}{j}\PY{+w}{ }\PY{o}{\PYZlt{}}\PY{+w}{ }\PY{n}{A}\PY{p}{.}\PY{n}{width}\PY{p}{;}\PY{+w}{ }\PY{n}{j}\PY{o}{+}\PY{o}{+}\PY{p}{)}
\PY{+w}{      }\PY{n}{A}\PY{p}{.}\PY{n}{elements}\PY{p}{[}\PY{n}{i}\PY{o}{*}\PY{n}{A}\PY{p}{.}\PY{n}{width}\PY{+w}{ }\PY{o}{+}\PY{+w}{ }\PY{n}{j}\PY{p}{]}\PY{+w}{ }\PY{o}{=}\PY{+w}{ }\PY{p}{(}\PY{k+kt}{float}\PY{p}{)}\PY{p}{(}\PY{n}{rand}\PY{p}{(}\PY{p}{)}\PY{+w}{ }\PY{o}{\PYZpc{}}\PY{+w}{ }\PY{l+m+mi}{3}\PY{p}{)}\PY{p}{;}

\PY{+w}{  }\PY{c+c1}{// Populate Matrix B with random floats (0.0 or 1.0)}
\PY{+w}{  }\PY{k}{for}\PY{p}{(}\PY{k+kt}{int}\PY{+w}{ }\PY{n}{i}\PY{+w}{ }\PY{o}{=}\PY{+w}{ }\PY{l+m+mi}{0}\PY{p}{;}\PY{+w}{ }\PY{n}{i}\PY{+w}{ }\PY{o}{\PYZlt{}}\PY{+w}{ }\PY{n}{B}\PY{p}{.}\PY{n}{height}\PY{p}{;}\PY{+w}{ }\PY{n}{i}\PY{o}{+}\PY{o}{+}\PY{p}{)}
\PY{+w}{    }\PY{k}{for}\PY{p}{(}\PY{k+kt}{int}\PY{+w}{ }\PY{n}{j}\PY{+w}{ }\PY{o}{=}\PY{+w}{ }\PY{l+m+mi}{0}\PY{p}{;}\PY{+w}{ }\PY{n}{j}\PY{+w}{ }\PY{o}{\PYZlt{}}\PY{+w}{ }\PY{n}{B}\PY{p}{.}\PY{n}{width}\PY{p}{;}\PY{+w}{ }\PY{n}{j}\PY{o}{+}\PY{o}{+}\PY{p}{)}
\PY{+w}{      }\PY{n}{B}\PY{p}{.}\PY{n}{elements}\PY{p}{[}\PY{n}{i}\PY{o}{*}\PY{n}{B}\PY{p}{.}\PY{n}{width}\PY{+w}{ }\PY{o}{+}\PY{+w}{ }\PY{n}{j}\PY{p}{]}\PY{+w}{ }\PY{o}{=}\PY{+w}{ }\PY{p}{(}\PY{k+kt}{float}\PY{p}{)}\PY{p}{(}\PY{n}{rand}\PY{p}{(}\PY{p}{)}\PY{+w}{ }\PY{o}{\PYZpc{}}\PY{+w}{ }\PY{l+m+mi}{2}\PY{p}{)}\PY{p}{;}

\PY{+w}{  }\PY{c+c1}{// Call the host function to execute the multiplication on the GPU}
\PY{+w}{  }\PY{n}{MatMul}\PY{p}{(}\PY{n}{A}\PY{p}{,}\PY{+w}{ }\PY{n}{B}\PY{p}{,}\PY{+w}{ }\PY{n}{C}\PY{p}{)}\PY{p}{;}

\PY{+w}{  }\PY{k}{return}\PY{+w}{ }\PY{l+m+mi}{0}\PY{p}{;}
\PY{p}{\PYZcb{}}

[Output]
> MatrixMultiplication_shared 1024 1024 1024
CUDA malloc A: no error
CUDA malloc B: no error
CUDA malloc C: no error
Run kernel: no error
Time: 16.10918 ms
Copy C off of device: no error
\end{Verbatim}